% !TeX program = xelatex
% 数学 LaTeX 排版 Skill v4.0.0；版式保持 v3 金标准. XeLaTeX 编译至少两次. 不分发字体.
% WZ-LOCKED-CLASS-BEGIN
\begin{filecontents*}[overwrite]{elegantbook-original-adapter.cls}
\NeedsTeXFormat{LaTeX2e}
\ProvidesClass{elegantbook-original-adapter}[2026/09/23 v3.0 fixed mathematical book environments]
\LoadClass[10pt,letterpaper,oneside]{book}
\RequirePackage[UTF8,scheme=plain,fontset=fandol]{ctex}
\RequirePackage{fontspec}
\setmainfont{cmunrm.otf}[BoldFont=cmunbx.otf,ItalicFont=cmunti.otf,BoldItalicFont=cmunbi.otf]
\setCJKfamilyfont{sourceheader}[ItalicFont=FandolKai-Regular.otf]{FandolKai-Regular.otf}
\RequirePackage{amsmath,amssymb,mathrsfs}
\RequirePackage{amsthm,etoolbox}
\RequirePackage{xcolor,geometry,setspace,fancyhdr,enumitem,needspace,xurl}
\definecolor{structurecolor}{HTML}{880E4F}
\definecolor{main}{HTML}{9C27B0}
\geometry{paperwidth=612bp,paperheight=792bp,left=20mm,right=20mm,top=95.04bp,bottom=20mm,headheight=12pt,headsep=25pt,footskip=30pt}
\setstretch{1.6}
\setlength{\parindent}{2em}
\setlength{\parskip}{0pt}
\setlist{nosep,leftmargin=2em}
\AtBeginDocument{\setlength{\abovedisplayskip}{4pt plus 1pt minus 1pt}\setlength{\belowdisplayskip}{4pt plus 1pt minus 1pt}\setlength{\abovedisplayshortskip}{2pt}\setlength{\belowdisplayshortskip}{3pt}}
\fancyhf{}
\fancyhead[L]{{\itshape\CJKfamily{sourceheader}\rightmark}}
\fancyhead[R]{{\itshape\CJKfamily{sourceheader}\leftmark}}
\fancyfoot[C]{\thepage}
\renewcommand{\headrulewidth}{0.4pt}
\renewcommand{\headrule}{{\color{structurecolor}\hrule height\headrulewidth width\headwidth\vskip-\headrulewidth}}
\pagestyle{fancy}
\fancypagestyle{cover}{\fancyhf{}\fancyfoot[C]{\thepage}\renewcommand{\headrulewidth}{0pt}\renewcommand{\headrule}{}}
\RequirePackage{hyperref}
\hypersetup{unicode=true,colorlinks=true,linkcolor=structurecolor,urlcolor=structurecolor,citecolor=structurecolor,linktoc=section,bookmarksnumbered=true,bookmarksopen=true,pdfborder={0 0 0},pdfstartview=FitH}
\setcounter{tocdepth}{1}
\renewcommand*\l@section{\@dottedtocline{1}{1.5em}{2.4em}}
\renewcommand{\thesection}{\thechapter.\arabic{section}}
\newcommand{\originalchapter}[1]{\par\addvspace{14pt}\Needspace{5\baselineskip}\refstepcounter{chapter}\setcounter{section}{0}\addcontentsline{toc}{chapter}{\protect\numberline{\thechapter}#1}\markboth{\thechapter\quad #1}{}\begingroup\centering\color{structurecolor}\bfseries\fontsize{14.4}{20}\selectfont\thechapter\quad #1\par\endgroup\nobreak\vspace{8pt}}
\newcommand{\originalsection}[1]{\par\addvspace{9pt}\Needspace{4\baselineskip}\refstepcounter{section}\addcontentsline{toc}{section}{\protect\hspace*{1.5em}\protect\numberline{\protect\textcolor{structurecolor}{\thesection}}#1}\markright{\thesection\quad #1}\noindent{\fontsize{12}{17}\selectfont\color{structurecolor}\thesection\quad}{\bfseries\fontsize{12}{17}\selectfont #1}\par\nobreak\vspace{4pt}}

% Semantic environments are registered once here, not re-designed in manuscripts.
\newtheoremstyle{wzstatement}{6pt}{5pt}
  {\normalfont\kaishu}{0pt}{\normalfont\bfseries\color{structurecolor}}
  {}{.7em}{\thmname{#1}\thmnumber{ #2}\thmnote{\normalfont\ (#3)}}
\newtheoremstyle{wzdefinition}{6pt}{5pt}
  {\normalfont\songti}{0pt}{\normalfont\bfseries\color{structurecolor}}
  {}{.7em}{\thmname{#1}\thmnumber{ #2}\thmnote{\normalfont\ (#3)}}
\newtheoremstyle{wzproblem}{7pt}{5pt}
  {\normalfont\songti}{2em}{\normalfont\bfseries\color{main}}
  {}{.7em}{\thmname{#1}\thmnumber{ #2}}
\newtheoremstyle{wzremark}{4pt}{4pt}
  {\normalfont\songti}{0pt}{\normalfont\bfseries}
  {}{.7em}{\thmname{#1}}
\theoremstyle{wzstatement}
\newtheorem{theorem}{定理}[chapter]
\newtheorem{lemma}[theorem]{引理}
\newtheorem{proposition}[theorem]{命题}
\newtheorem{corollary}[theorem]{推论}
\theoremstyle{wzdefinition}
\newtheorem{definition}[theorem]{定义}
\theoremstyle{wzproblem}
\newtheorem{example}{例题}[chapter]
\newtheorem{exercise}{习题}[chapter]
\theoremstyle{wzremark}
\newtheorem*{remark}{注}
\renewcommand{\proofname}{证明}
% Retain the native amsthm QED stack; only adjust the fixed typography.
\expandafter\patchcmd\csname\string\proof\endcsname{\normalfont \topsep6\p@\@plus6\p@\relax}
  {\normalfont\songti \topsep5\p@\@plus1\p@\relax}
  {}{\ClassError{elegantbook-original-adapter}{Cannot patch proof spacing}{Check the amsthm version.}}
\expandafter\patchcmd\csname\string\proof\endcsname{\itshape}{\normalfont\bfseries}
  {}{\ClassError{elegantbook-original-adapter}{Cannot patch proof heading}{Check the amsthm version.}}
\expandafter\patchcmd\csname\string\proof\endcsname{\ignorespaces}
  {\setlength{\parindent}{2em}\ignorespaces}
  {}{\ClassError{elegantbook-original-adapter}{Cannot patch proof paragraphs}{Check the amsthm version.}}
\AtBeginEnvironment{proof}{\par\Needspace{3\baselineskip}}
\renewcommand{\qedsymbol}{\ensuremath{\square}}
\newenvironment{solution}{\begin{proof}[解答]}{\end{proof}}
\newcommand{\wzcheckchapter}{\ifnum\value{chapter}<1
  \ClassError{elegantbook-original-adapter}{Numbered mathematics before chapter 1}{Move it after the first originalchapter.}\fi}

\newcommand{\wzregisterguard}[1]{\AtBeginEnvironment{#1}{\wzcheckchapter\par\Needspace{3\baselineskip}}}
\forcsvlist{\wzregisterguard}{theorem,lemma,proposition,corollary,definition,example,exercise}
% Front matter and bibliography are not numbered mathematical chapters.
\newcommand{\originalpreface}{\par\addvspace{10pt}\Needspace{5\baselineskip}
  \noindent{\bfseries\fontsize{12}{17}\selectfont 前言}\par\nobreak\vspace{4pt}}
\newcommand{\originalcontents}{\par\addvspace{14pt}
  \begin{center}{\color{structurecolor}\bfseries\fontsize{14.4}{20}\selectfont 目录}\end{center}
  \vspace{-10pt}\@starttoc{toc}}
\renewcommand{\bibname}{参考文献}
\renewenvironment{thebibliography}[1]{
  \par\addvspace{12pt}\Needspace{3\baselineskip}\phantomsection
  \addcontentsline{toc}{chapter}{\bibname}
  \markboth{\bibname}{}
  \noindent{\color{structurecolor}\bfseries\fontsize{14.4}{20}\selectfont\bibname}\par\nobreak\vspace{6pt}
  \list{\@biblabel{\@arabic\c@enumiv}}
    {\settowidth\labelwidth{\@biblabel{#1}}\leftmargin\labelwidth
     \advance\leftmargin\labelsep\usecounter{enumiv}
     \let\p@enumiv\@empty\renewcommand\theenumiv{\@arabic\c@enumiv}
     \setlength{\itemsep}{3pt}\setlength{\parsep}{0pt}}
  \sloppy\clubpenalty4000\widowpenalty4000\sfcode`\.=1000\relax
}{\def\@noitemerr{\@latex@warning{Empty bibliography}}\endlist}

\numberwithin{equation}{chapter}
\renewcommand{\theequation}{\thechapter.\arabic{equation}}
\allowdisplaybreaks[2]
\emergencystretch=2em
\clubpenalty=6000
\widowpenalty=6000
\raggedbottom
\endinput
\end{filecontents*}
% WZ-LOCKED-CLASS-END
\documentclass{elegantbook-original-adapter}
% ===== 元数据内容槽 =====
\newcommand{\BookTitle}{抽象代数: 群、环、模与域}
\hypersetup{pdftitle={抽象代数: MIT公开讲义中文重构本},pdfauthor={James McKernan; MIT OCW学生笔记; 中文重构}}
\usepackage{tikz,booktabs,longtable}
\usetikzlibrary{arrows.meta,positioning,calc}
\newcommand{\Z}{\mathbb Z}
\newcommand{\Q}{\mathbb Q}
\newcommand{\R}{\mathbb R}
\newcommand{\C}{\mathbb C}
\newcommand{\F}{\mathbb F}
\newcommand{\N}{\mathbb N}
\newcommand{\Aut}{\operatorname{Aut}}
\newcommand{\Gal}{\operatorname{Gal}}
\newcommand{\Hom}{\operatorname{Hom}}
\newcommand{\im}{\operatorname{im}}
\newcommand{\ord}{\operatorname{ord}}
\newcommand{\sgn}{\operatorname{sgn}}
\newcommand{\lcm}{\operatorname{lcm}}
\newcommand{\ann}{\operatorname{Ann}}
\begin{document}
\frontmatter
\begin{center}
{\color{structurecolor}\bfseries\fontsize{18}{25}\selectfont\BookTitle\par}
{\bfseries MIT公开讲义中文重构本\par}
James McKernan 原课程讲义; Jakin Ng、Sanjana Das、Ethan Yang 学生笔记补充\par
18.703 (2013), RES.18-011 (2021), RES.18-012 (2022)\par
\end{center}
\originalpreface
群记录可逆运算, 环把加法和乘法放在同一结构中, 模把向量空间的标量域放宽为环, 域扩张则把多项式的根纳入可研究的空间. 本书按这条依赖关系组织本科抽象代数. 读者需要整数整除、初等多项式及线性代数中的基和维数; 所用域上的线性代数论证在使用处重述. 不要求先学表示论、交换代数或代数数论.

主源为 James McKernan 的 MIT OpenCourseWare 18.703 Modern Algebra, Spring 2013 的23份教师讲义. 群作用移到 Sylow 定理之前; 商对象的共同逻辑通过群、环、模逐次展开. 补充来源为 RES.18-011 Algebra I Student Notes, Fall 2021 和 RES.18-012 Algebra II Student Notes, Spring 2022 的合订学生笔记及官方 TeX 包. 官网明确说明两组学生笔记未经任课教师或其他 MIT 教师核准. 这里按数学命题重新组织、补足证明并独立交叉审读; 不把其刊载身份当作正确性保证. 英文原件只在仓库 sources/abstract-algebra 中归档, 中文源码包不含英文原讲义、截图或原 TeX.

本书完整论证群同态与商群、群作用与 Sylow 定理、环与理想、唯一分解与多项式、主理想整环上的有限生成模、域扩张、有限域以及有限 Galois 对应. 自设例题和带解答练习用于复核定义及推理; 与原讲义相同的经典数学事实不宣称为原创. 不重写数值线性代数, 不展开特征标理论. 原学生笔记中的欧氏几何、Lie 群、谱定理、代数数论的理想类群、函数域几何以及一般根式可解性定理不在本册范围. 书末按原课逐讲标明完整覆盖、择要覆盖和未覆盖内容; 未覆盖部分不以几句名称代替证明.

资料来自 Massachusetts Institute of Technology 的 MIT OpenCourseWare, 默认许可为 \href{https://creativecommons.org/licenses/by-nc-sa/4.0/}{CC BY-NC-SA 4.0}. 中文重构、补充证明、另设例题及原生 TikZ 图以相同许可作非商业使用和相同方式共享, 不表示 MIT、教师或学生记录者认可本项目. 单独第三方声明优先; 不复用外书题面或受限图片. 署名与逐文件许可核查见随附源清单和 LICENSE.md; 官方条款见 \href{https://ocw.mit.edu/pages/privacy-and-terms-of-use/}{MIT OCW 使用条款}. 版本日期: 2026-10-08.
\clearpage
\phantomsection\addcontentsline{toc}{chapter}{课程资料与教学进度}
\noindent{\color{structurecolor}\bfseries\fontsize{14.4}{20}\selectfont 课程资料与教学进度}\par
主课程为 MIT 数学系本科课程18.703“现代代数”, James McKernan 教授, 2013年春季学期. 官方安排每周两次讲课, 每次1.5小时. 官网先修列为18.02多元微积分; 18.700线性代数与18.703共同组成标准代数序列. 本中文册在模和域部分使用线性代数基础, 这是补充来源带来的阅读要求, 不改写成主课程官方先修.

大纲为有限群理论, 环、理想、欧几里得环和多项式的唯一分解, 域及有限域的性质和应用. 指定教材为 I. N. Herstein, \textit{Abstract Algebra}, Macmillan, 1986, ISBN 9780023538209. 阅读页另列 Tom Judson 的在线教材 \textit{Abstract Algebra: Theory and Applications} 各节作为课前阅读. 仅有书目题号的作业不代表官网公开了相应教材题面.

考核为作业50\%, 两次期中各10\%, 期末30\%. 课程说明作业通常周二布置、下一周周四交, 考试周可例外; 不收迟交, 去掉最低一次作业分. 考试有复述定义部分. 以下按照官方教学进度表的讲次和 E1/E2/E3 考试标签列出, 官网未列具体日历日期. 原作业文件的截止日期另随题集登记, 不用它反推每次讲课日期.
\begin{longtable}{@{}p{.09\textwidth}p{.45\textwidth}p{.16\textwidth}p{.22\textwidth}@{}}
\toprule 讲次&主题&交作业&Judson阅读节\\\midrule\endhead
1&群&&3.1,3.2\\
2&子群&&3.3\\
3&陪集&&6.1,6.2\\
4&循环群&作业1&6.2,4.1\\
5&置换群&&4.2,5.1\\
6&对称群中的共轭&&5.1,5.2\\
7&同构&作业2&9.1\\
8&同态与核&&11.1,10.1\\
9&商群&作业3&10.1,11.2\\
10&同构定理&&11.2\\
11&交错群&作业4&9.2\\
12&表示与小阶群&&13.1\\
13&Sylow定理与应用&作业5&15.1,15.2\\
14&环&作业6&16.1,16.2\\
E1&第一次期中&&\\
15&环的基本性质&&期中复习 (阅读页原列)\\
16&环同态与理想&作业7&16.3\\
17&分式域&作业8&18.1,16.4\\
18&素理想与极大理想&&18.2\\
19&特殊整环&作业9&18.2,17.2\\
20&欧几里得整环&&17.2,17.3\\
E2&第二次期中&作业10&\\
21&多项式环&&17.2,17.3\\
22&快速素性检验&作业11&17.3\\
23&群作用与自同构&&期中复习 (阅读页原列)\\
24&复习, 无讲义&&17.3\\
E3&期末考试&&\\
\bottomrule
\end{longtable}
补充资料 RES.18-011 是2021年秋季18.701“代数I”的学生笔记, 讲课教师 Davesh Maulik, 记录者 Jakin Ng、Sanjana Das、Ethan Yang, 合订178页. RES.18-012 是2022年春季18.702“代数II”的学生笔记, 讲课教师 Roman Bezrukavnikov, 记录者 Sanjana Das、Jakin Ng, 笔记工作由 Ashay Athalye 监督, 合订162页. 两个资源页均明确未经教师核准; 此处“监督记录”不等于审定数学内容. 两册各35讲, 第二册另有特征标附录. 没有在这两组资源的页面公开独立作业或考试卷; 不把其所链接的其他年份课程材料冒充2021/2022年的试题. 本册只按明确选定的补充讲次使用.

参考: \href{https://ocw.mit.edu/courses/18-703-modern-algebra-spring-2013/pages/syllabus/}{官方课程大纲}, \href{https://ocw.mit.edu/courses/18-703-modern-algebra-spring-2013/pages/calendar/}{教学进度}, \href{https://ocw.mit.edu/courses/18-703-modern-algebra-spring-2013/pages/readings/}{阅读安排}.

\clearpage
\originalcontents
\clearpage
\mainmatter
\originalchapter{群与子群}
\originalsection{可逆运算的公理}
正三角形的旋转和翻折可以复合, 且每一步可以撤销. 把这类运算抽象出来, 需要先约定复合次序: 本书 $gh$ 表示先作用 $h$, 再作用 $g$. 加法群中相应运算写作 $a+b$.
\begin{definition}
群是带二元运算的非空集合 $G$, 满足结合律 $(ab)c=a(bc)$, 存在单位元 $e$ 使 $ea=ae=a$, 且每个 $a$ 有逆元 $a^{-1}$ 使 $aa^{-1}=a^{-1}a=e$. 若还满足 $ab=ba$, 则称为交换群或阿贝尔群.
\end{definition}
\begin{proposition}
单位元和每个元素的逆元唯一; $(ab)^{-1}=b^{-1}a^{-1}$. 群中左右消去律成立.
\end{proposition}
\begin{proof}
若 $e,f$ 都是单位元, 则 $e=ef=f$. 若 $b,c$ 都是 $a$ 的逆元, 则 $b=b(ac)=(ba)c=c$. 两个乘积 $(ab)(b^{-1}a^{-1})$ 与 $(b^{-1}a^{-1})(ab)$ 都为 $e$, 给出乘积逆元公式. 对 $ax=ay$ 左乘 $a^{-1}$, 对 $xa=ya$ 右乘 $a^{-1}$, 分别得 $x=y$.
\end{proof}
整数、有理数、实数与复数在加法下均成群. 非零有理数在乘法下成群, 整数在乘法下不成群: $2$ 没有整数逆元. 域 $F$ 上的可逆 $n$ 阶矩阵成群 $\mathrm{GL}_n(F)$; 矩阵结合律保证复合结合, 可逆性保证乘积仍在集合内. 一般矩阵乘法不交换.
\begin{definition}
非空子集 $H\subseteq G$ 若在限制运算下成群, 称为子群, 记 $H\le G$. 群的阶 $|G|$ 是其元素数, 可为无穷.
\end{definition}
\begin{proposition}[子群判别]
非空子集 $H$ 是子群当且仅当对所有 $a,b\in H$, 有 $ab^{-1}\in H$.
\end{proposition}
\begin{proof}
必要性由子群封闭性得到. 反之取 $a\in H$, 得 $e=aa^{-1}\in H$; 再令第一元素为 $e$ 得 $b^{-1}\in H$; 最后对 $a,b^{-1}$ 应用条件得 $ab\in H$. 结合律继承自 $G$, 单位元和逆元也都属于 $H$.
\end{proof}
\originalsection{生成与循环群}
任意一族子群的交仍是子群: 单位元在每个子群中, 且判别式 $ab^{-1}$ 同时留在每个子群中. 包含集合 $S$ 的所有子群的交称为 $\langle S\rangle$. 它恰由 $S$ 中元素及其逆元的有限乘积组成, 包括空乘积 $e$: 这些乘积成子群, 又被任何包含 $S$ 的子群包含.
\begin{definition}
由一个元素 $a$ 生成的群 $\langle a\rangle=\{a^k:k\in\Z\}$ 称为循环群. 元素 $a$ 的阶是使 $a^m=e$ 的最小正整数 $m$; 若不存在则称阶无穷.
\end{definition}
\begin{theorem}\label{thm:cyclic}
整数加法群的子群均为 $d\Z$, 其中 $d\ge0$. 无限循环群同构于 $\Z$, 阶为 $n$ 的循环群同构于 $\Z/n\Z$. 循环群的子群仍循环; 有限循环群中对每个 $d\mid n$ 恰有一个阶为 $d$ 的子群.
\end{theorem}
\begin{proof}
非零子群 $H\le\Z$ 含正整数. 取最小正元素 $d$, 对任何 $h\in H$ 作除法 $h=qd+r$, $0\le r<d$. 因 $r\in H$, 最小性迫使 $r=0$, 故 $H=d\Z$. 零子群对应 $d=0$.

映射 $k\mapsto a^k$ 满射到 $\langle a\rangle$. 若 $a$ 阶无限则它单射; 若阶为 $n$, 除法证明 $a^k=e$ 当且仅当 $n\mid k$, 故余数类恰对应不同元素. 对任意子群 $K\le\langle a\rangle$, 原像是 $d\Z$, 因而 $K=\langle a^d\rangle$. 有限情形原像包含 $n\Z$, 所以 $d\mid n$, 且 $K$ 有 $n/d$ 个元素. 给定所需阶 $r\mid n$, 必须且只需取 $d=n/r$.
\end{proof}
\begin{example}
求 $\Z/18\Z$ 中 $\overline{12}$ 的阶及其生成的子群.
\end{example}
\begin{solution}
$k\overline{12}=0$ 等价于 $18\mid12k$, 即 $3\mid2k$, 即 $3\mid k$. 阶为3, 子群为 $\{\bar0,\bar6,\overline{12}\}$. 一般 $\bar a$ 在 $\Z/n\Z$ 的阶为 $n/\gcd(a,n)$: 约去最大公因子后两整数互素.
\end{solution}
\originalsection{带解答练习}
\begin{exercise}证明有限群中非空、乘法封闭的子集是子群. 说明无限情形为何不能照搬.\end{exercise}
\begin{solution}
对 $a\in H$, 正幂序列中两项相等, 设 $a^i=a^j$, $i<j$. 在 $G$ 内消去得 $a^{j-i}=e\in H$. 若 $a\ne e$, 则 $a^{-1}=a^{j-i-1}\in H$; $a=e$ 时逆元本身已在 $H$. 于是逆元和乘积封闭. 无限群 $\Z$ 的正整数子集加法封闭, 却没有0和加法逆元.
\end{solution}
\begin{exercise}设 $a$ 的有限阶为 $n$. 证明 $\ord(a^k)=n/\gcd(n,k)$.\end{exercise}
\begin{solution}
写 $n=dn'$, $k=dk'$, $\gcd(n',k')=1$. 条件 $(a^k)^m=e$ 等价于 $n\mid km$, 等价于 $n'\mid k'm$, 再由互素性等价于 $n'\mid m$. 最小正 $m$ 为 $n'$.
\end{solution}

\originalchapter{陪集、同态与商群}
\originalsection{陪集与拉格朗日定理}
\begin{definition}
对 $H\le G$, 左陪集 $gH=\{gh:h\in H\}$, 右陪集 $Hg=\{hg:h\in H\}$. 左陪集数称为指数 $[G:H]$.
\end{definition}
\begin{lemma}
左陪集或者相等或者不交; $aH=bH$ 当且仅当 $b^{-1}a\in H$. 每个陪集与 $H$ 等势.
\end{lemma}
\begin{proof}
若 $ah_1=bh_2$, 则 $b^{-1}a=h_2h_1^{-1}\in H$, 所以 $a=bh$, 从而 $aH=bH$. 反之 $b^{-1}a\in H$ 即给出相等. 映射 $h\mapsto ah$ 是双射, 逆映射为 $x\mapsto a^{-1}x$. 每个元素在自己的陪集中, 故这些不交陪集划分 $G$.
\end{proof}
\begin{theorem}[拉格朗日]
有限群中 $|G|=[G:H]|H|$. 特别地, 子群阶和元素阶都整除群阶.
\end{theorem}
\begin{proof}
对陪集划分计数得第一个等式. 再对循环子群 $\langle a\rangle$ 应用即可.
\end{proof}
逆命题不能任意使用: 群阶的一个因子未必是某个子群的阶. 后文将用 $A_4$ 给出阶6的反例, 用 Sylow 定理说明素数幂最大因子的存在.
\originalsection{同态与正规子群}
\begin{definition}
映射 $\varphi:G\to G'$ 若满足 $\varphi(ab)=\varphi(a)\varphi(b)$, 称为同态; 双射同态称为同构. 核为 $\ker\varphi=\{g:\varphi(g)=e\}$, 像为 $\im\varphi=\varphi(G)$. 子群 $N$ 若对每个 $g\in G$ 有 $gNg^{-1}=N$, 称为正规子群, 记 $N\trianglelefteq G$.
\end{definition}
同态把单位元送到单位元: $\varphi(e)=\varphi(e)^2$, 消去即可. 又把逆元送到逆元. 因此像和核均是子群; 对 $n\in\ker\varphi$, $\varphi(gng^{-1})=e$, 所以核正规. 同态单射等价于核只有单位元: 两个元素像相等当且仅当 $b^{-1}a\in\ker\varphi$.
\begin{theorem}[商群的构造]\label{thm:quotient}
若 $N\trianglelefteq G$, 则陪集集合 $G/N$ 在 $(aN)(bN)=abN$ 下成群; 投影 $\pi:g\mapsto gN$ 的核为 $N$. 反之, 此乘法若良定义, 则 $N$ 必正规.
\end{theorem}
\begin{proof}
改变代表元为 $an_1,bn_2$ 时, 乘积是 $ab(b^{-1}n_1b)n_2$, 与 $ab$ 属同一陪集, 这里恰用到正规性. 结合律来自 $G$, 单位元为 $N$, 逆元为 $a^{-1}N$. 投影保持乘积且 $gN=N$ 当且仅当 $g\in N$.

反之 $(gN)(N)(g^{-1}N)=N$ 必须不依赖 $N$ 的代表元, 所以任何 $gng^{-1}$ 都属于 $N$. 这给出 $gNg^{-1}\subseteq N$; 用 $g^{-1}$ 再作一次得到反向包含.
\end{proof}
\begin{theorem}[第一同构定理]\label{thm:iso1}
任何群同态诱导同构 $G/\ker\varphi\cong\im\varphi$, 对应 $g\ker\varphi\mapsto\varphi(g)$.
\end{theorem}
\begin{proof}
陪集相同意味着 $h^{-1}g$ 在核中, 故像相同; 反之像相同也给出陪集相同. 因而诱导映射良定义且单射, 按像的定义满射, 并保持乘积.
\end{proof}
\originalsection{同构定理与对应关系}
\begin{theorem}
设 $N\trianglelefteq G$, $H\le G$. 则 $HN=\{hn:h\in H,n\in N\}$ 是子群, $H\cap N\trianglelefteq H$, 且 $H/(H\cap N)\cong HN/N$. 若 $N\le K\trianglelefteq G$, 则 $(G/N)/(K/N)\cong G/K$. 包含 $N$ 的子群 $H$ 与 $G/N$ 的子群 $\bar H$ 一一对应, 对应为 $H\mapsto H/N$, $\bar H\mapsto\pi^{-1}(\bar H)$; 正规性也对应.
\end{theorem}
\begin{proof}
两个 $hn$ 的乘积可用 $n_1h_2=h_2(h_2^{-1}n_1h_2)$ 改写为同一形式, 逆元也可改写, 所以 $HN$ 成群. 限制投影 $H\to HN/N$ 满射且核为 $H\cap N$, 用第一同构定理. 映射 $G/N\to G/K$, $gN\mapsto gK$, 良定义、满射且核为 $K/N$, 再用第一同构定理.

子群原像总成子群. 若 $H$ 包含 $N$, 则 $\pi^{-1}(\pi(H))=H$, 因 $gN=hN$ 意味着 $g\in hN\subseteq H$. 满射性又给 $\pi(\pi^{-1}(\bar H))=\bar H$. 最后 $\pi(gHg^{-1})=\pi(g)\bar H\pi(g)^{-1}$, 结合原像相等式即得正规性对应.
\end{proof}
\begin{example}求 $\Z/12\Z$ 到 $\Z/8\Z$ 的所有群同态.\end{example}
\begin{solution}
同态由 $\bar1$ 的像 $x$ 决定, 且必须满足 $12x=0$ 模8. 这等价于 $4x=0$, 即 $x$ 为偶数. 四个同态分别由 $x=0,2,4,6$ 给出. 定义 $\bar k\mapsto kx$ 不依赖整数代表元, 因为 $12x=0$.
\end{solution}
\originalsection{带解答练习}
\begin{exercise}证明指数为2的子群正规.\end{exercise}
\begin{solution}
若 $g\in H$, 左右陪集均为 $H$. 若 $g\notin H$, 左陪集 $gH$ 与右陪集 $Hg$ 都是 $G\setminus H$: 任一方向的陪集划分都只有两块. 因而 $gH=Hg$, 等价于正规性.
\end{solution}
\begin{exercise}若 $H,K\trianglelefteq G$ 且 $H\cap K=\{e\}$, 证明 $hk=kh$.\end{exercise}
\begin{solution}
交换子 $hkh^{-1}k^{-1}$ 因 $K$ 正规属于 $K$, 又因 $H$ 正规可写成 $h(kh^{-1}k^{-1})\in H$. 交只有 $e$, 所以交换子为 $e$, 即所求.
\end{solution}

\originalchapter{置换群与群的具体构造}
\originalsection{循环分解与符号}
有限集合 $\{1,\ldots,n\}$ 上的双射组成对称群 $S_n$, 阶为 $n!$. 对任一点反复作用置换, 首次重复必返回起点; 各轨道互不相交, 于是置换唯一分解为不交循环, 只差循环书写起点及各循环次序. 不交循环交换, 一个长度 $r$ 的循环的阶是 $r$, 所以整个置换的阶为各循环长度的最小公倍数.

长度 $r$ 的循环可以写为 $(a_1\,a_r)\cdots(a_1\,a_2)$, 即 $r-1$ 个换位. 要定义奇偶性, 还必须证明不同换位分解的奇偶一致.
\begin{proposition}
置换矩阵 $P_\sigma$, 由 $P_\sigma e_i=e_{\sigma(i)}$ 定义, 满足 $P_{\sigma\tau}=P_\sigma P_\tau$. 因而 $\sgn(\sigma)=\det P_\sigma\in\{1,-1\}$ 是同态, 每个换位符号为 $-1$. 换位分解的奇偶性唯一.
\end{proposition}
\begin{proof}
在每个基向量上检查复合等式. 行列式乘法公式给出同态性. 置换矩阵由单位矩阵交换列得到, 行列式只能为 $\pm1$; 单次换位交换两列, 行列式为 $-1$. 所以任意 $k$ 个换位的乘积行列式为 $(-1)^k$, 不依赖分解. 这里仅使用先修线性代数的行列式性质.
\end{proof}
\begin{definition}
交错群 $A_n=\ker\sgn$ 由偶置换构成. $n\ge2$ 时符号满射, 故 $|A_n|=n!/2$, 且 $A_n\trianglelefteq S_n$.
\end{definition}
\originalsection{共轭与凯莱定理}
\begin{proposition}
$\tau(a_1\,\cdots\,a_r)\tau^{-1}=(\tau(a_1)\,\cdots\,\tau(a_r))$. 两个置换在 $S_n$ 中共轭当且仅当其循环长度多重集相同.
\end{proposition}
\begin{proof}
右侧把 $\tau(a_i)$ 送至 $\tau(a_{i+1})$, 与左侧一致, 其余点均固定. 因而共轭保留循环类型. 反之把长度对应循环的点依次一一配对, 得置换 $\tau$, 上述公式即给出共轭.
\end{proof}
\begin{theorem}[凯莱]
每个群同构于某个置换群的子群; 有限阶 $n$ 的群可嵌入 $S_n$.
\end{theorem}
\begin{proof}
对 $g\in G$ 定义左平移 $L_g(x)=gx$. 逆为 $L_{g^{-1}}$, 故是置换. $L_gL_h=L_{gh}$, 所以 $g\mapsto L_g$ 是同态; 若 $L_g=L_h$, 在 $e$ 上取值得 $g=h$. 有限时将 $G$ 的元素编号即可.
\end{proof}
\originalsection{直积、二面体群与表示}
群 $G,H$ 的直积 $G\times H$ 逐坐标乘法. $\Z/m\Z\times\Z/n\Z$ 循环当且仅当 $\gcd(m,n)=1$: 元素 $(\bar1,\bar1)$ 的阶是 $\lcm(m,n)$; 若不互素, 任一元素阶都至多该最小公倍数, 小于群阶.

正 $n$ 边形 ($n\ge3$) 的旋转 $r$ 与一个翻折 $s$ 满足 $r^n=s^2=e$, $srs=r^{-1}$. 任何词用 $sr=r^{-1}s$ 可移成 $r^i$ 或 $r^is$, 至多 $2n$ 种. 几何上这 $2n$ 个变换不同, 组成二面体群 $D_n$, 本书约定其阶为 $2n$. 此例也解释生成元和关系的表示 $\langle r,s\mid r^n=s^2=e,srs=r^{-1}\rangle$: 把所有形式词对关系造成的等同取商. 一般自由群构造不在此详述, 后文只在已验证的具体群中使用表示.
\begin{example}确定 $S_4$ 的共轭类和 $A_4$ 的一个正规子群, 并证明 $A_4$ 没有阶6子群.\end{example}
\begin{solution}
循环类型为 $1^4,2\,1^2,2^2,3\,1,4$, 元素数分别为 $1,6,3,8,6$, 总计24. $V=\{e,(12)(34),(13)(24),(14)(23)\}$ 是子群; 共轭保留双换位, 故正规, 并且 $V\cong C_2\times C_2$.

若 $H\le A_4$ 阶6, 则指数2, 正规. 商 $A_4/H$ 阶2, 所有阶3元素在商中的像为单位元, 因为像的阶同时整除3与2. 但 $A_4$ 的八个3循环及单位元都会在 $H$ 中, 超过6, 矛盾.
\end{solution}
\originalsection{带解答练习}
\begin{exercise}求 $(1234)(56)$ 的阶与符号, 并列出 $S_3$ 的全部正规子群.\end{exercise}
\begin{solution}
阶为 $\lcm(4,2)=4$, 符号为 $(-1)^3(-1)=1$. $S_3$ 子群阶仅可为1,2,3,6. 阶2子群由一个换位生成, 被其他换位共轭到不同子群, 不正规. 唯一阶3子群是 $A_3$, 正规. 故正规子群为 $\{e\},A_3,S_3$.
\end{solution}
\begin{exercise}证明 $D_n$ 的旋转子群正规, 并计算中心.\end{exercise}
\begin{solution}
旋转子群指数2. $r^i$ 与 $s$ 交换等价于 $r^{2i}=e$. 而 $r^is$ 与 $r$ 交换要求 $r=r^{-1}$, 对 $n\ge3$ 不成立. 因而 $n$ 奇数时中心只有 $e$; $n$ 偶数时中心为 $\{e,r^{n/2}\}$.
\end{solution}

\originalchapter{群作用与 Sylow 定理}
\originalsection{轨道与稳定子}
\begin{definition}
若有限群阶为素数 $p$ 的幂, 称为 $p$ 群; $p$ 子群是自身为 $p$ 群的子群. 中心 $Z(G)$ 是与群中每个元素交换的元素集合; 中心化子 $C_G(x)$ 是与指定元素 $x$ 交换的元素集合, 两者由子群判别法成子群. 自同构是从群到自身的同构, 所有自同构在复合下成群 $\Aut(G)$. 非平凡群若没有非平凡真正规子群, 称单群.
\end{definition}
\begin{definition}
群 $G$ 在集合 $X$ 上的作用是规则 $(g,x)\mapsto gx$, 满足 $ex=x$ 和 $(gh)x=g(hx)$. 点 $x$ 的轨道是 $Gx$, 稳定子是 $G_x=\{g:gx=x\}$.
\end{definition}
稳定子由子群判别法成群. “在同一轨道”是等价关系: 单位元、逆元和乘积分别给出自反、对称和传递性, 故轨道划分 $X$.
\begin{theorem}[轨道稳定子公式]\label{thm:orbit}
映射 $G/G_x\to Gx$, $gG_x\mapsto gx$, 是双射. 有限情形 $|Gx|=[G:G_x]$.
\end{theorem}
\begin{proof}
$gx=hx$ 当且仅当 $h^{-1}g\in G_x$, 恰是陪集相等条件; 因而映射良定义、单射且满射.
\end{proof}
共轭作用 $g\cdot x=gxg^{-1}$ 的稳定子是中心化子 $C_G(x)$, 单点轨道恰是中心 $Z(G)$ 的元素. 于是有限群满足类方程
\[
|G|=|Z(G)|+\sum_x[G:C_G(x)],
\]
其中求和对非中心共轭类各取一个代表.
\begin{lemma}\label{lem:pfix}
有限 $p$ 群作用于有限集合 $X$, 则 $|X|\equiv|X^G|\pmod p$, 其中 $X^G$ 为所有群元素都固定的点. 非平凡有限 $p$ 群的中心非平凡.
\end{lemma}
\begin{proof}
轨道大小整除群阶, 所以非单点轨道大小被 $p$ 整除; 把轨道大小相加即得同余. 对共轭作用, 固定点集为中心. 若 $p\mid|G|$, 则 $p\mid|Z(G)|$, 而中心包含 $e$, 故至少有 $p$ 个元素.
\end{proof}
\originalsection{Sylow 定理的证明}
\begin{theorem}[柯西]\label{thm:cauchy}
若素数 $p$ 整除有限群阶, 则群中存在阶为 $p$ 的元素.
\end{theorem}
\begin{proof}
令 $X=\{(g_1,\ldots,g_p):g_1\cdots g_p=e\}$. 前 $p-1$ 项任取, 最后一项唯一, 所以 $|X|=|G|^{p-1}$ 被 $p$ 整除. 循环移位保持乘积为 $e$: 移位后乘积为 $g_1^{-1}(g_1\cdots g_p)g_1=e$. 因此阶 $p$ 的循环群作用于 $X$. 固定点为 $(g,\ldots,g)$ 且 $g^p=e$. 引理\ref{lem:pfix}说明这类 $g$ 的个数被 $p$ 整除, 至少含 $e$, 故另有非单位元, 其阶恰为 $p$.
\end{proof}
\begin{theorem}[Sylow]\label{thm:sylow}
设 $|G|=p^am$, $p\nmid m$. 则有阶 $p^a$ 的子群, 称 Sylow $p$ 子群. 任一 $p$ 子群包含于某个 Sylow $p$ 子群, 所有 Sylow $p$ 子群互相共轭. 其个数 $n_p$ 满足 $n_p\mid m$ 且 $n_p\equiv1\pmod p$.
\end{theorem}
\begin{proof}
先归纳群阶证明存在性. 若 $a=0$, 单位子群即可. 若 $p\mid|Z(G)|$, 柯西定理给出中心内阶 $p$ 子群 $C$. 商 $G/C$ 的阶为 $p^{a-1}m$, 归纳得到其中阶 $p^{a-1}$ 子群, 原像阶为 $p^a$.

若 $p\nmid|Z(G)|$, 类方程至少有一项 $[G:C_G(x)]$ 不被 $p$ 整除, 否则右侧模 $p$ 不为0. 对应的中心化子是真子群, 其阶仍被 $p^a$ 整除, 在该真子群中归纳可得所求.

固定 Sylow 子群 $P$, 让任意 $p$ 子群 $Q$ 左乘作用于 $G/P$. 该集合有 $m$ 个点, 不被 $p$ 整除, 故有固定陪集 $gP$. 固定条件是 $g^{-1}Qg\le P$, 所以 $Q\le gPg^{-1}$. 若 $Q$ 本身也是 Sylow 子群, 阶相等迫使等号, 共轭性成立.

$G$ 共轭作用于所有 Sylow 子群, 轨道就是整个集合, 稳定子是正规化子 $N_G(P)=\{g:gPg^{-1}=P\}$. 因 $P\le N_G(P)$, $n_p=[G:N_G(P)]$ 整除 $m$. 再让 $P$ 作用于此集合. 若 $Q$ 被 $P$ 固定, 则 $P\le N_G(Q)$. 此时 $PQ$ 是子群, 因 $P$ 正规化 $Q$; 其阶 $|P||Q|/|P\cap Q|$ 是 $p$ 幂, 且不超过 $p^a$. 因 $Q$ 已阶 $p^a$, 必有 $P=Q$. 所以只有 $P$ 自己为固定点, 引理\ref{lem:pfix}得 $n_p\equiv1\pmod p$.
\end{proof}
\originalsection{小阶群与应用}
Sylow 子群唯一当且仅当正规, 因全部共轭. 若 $|G|=pq$, $p<q$ 为素数, 则 $n_q\mid p$ 且 $n_q\equiv1\pmod q$, 只能是1, 因而 $G$ 不是单群. 若还 $p\nmid q-1$, 阶 $q$ 正规子群 $Q$ 上的共轭作用给出 $P\to\Aut(Q)$, 而 $\Aut(C_q)$ 阶 $q-1$ (自同构由一个生成元的像决定, 该像可为任一非单位元, 共 $q-1$ 个; 每种选择给出可逆同态), 所以像阶只能为1. 于是 $P,Q$ 交换且交平凡, $G\cong P\times Q\cong C_{pq}$.
\begin{example}分类阶为 $p^2$ 的群.\end{example}
\begin{solution}
中心非平凡. 若中心阶 $p$, 则商群阶 $p$, 循环. 当 $G/Z(G)$ 循环时, 任意元素可写 $g^iz$, 两个这样元素交换, 所以 $G$ 交换, 反而中心就是 $G$. 因此中心不能只阶 $p$, 群交换. 若有阶 $p^2$ 元素则循环; 否则所有非单位元阶 $p$. 取 $a\ne e$, 再取 $b\notin\langle a\rangle$, 两个阶 $p$ 子群交只有 $e$, 其直积给出整个 $G$, 故 $G\cong C_p\times C_p$.
\end{solution}
\originalsection{带解答练习}
\begin{exercise}证明阶15的群循环; 证明阶12的群必有正规 Sylow 子群或正规阶4子群.\end{exercise}
\begin{solution}
$15=3\cdot5$, 且 $3\nmid4$, 用前述结论. 对阶12, $n_3=1$ 或4. 若为1已有正规子群; 若为4, 各阶3子群仅交于 $e$, 共给出8个阶3元素. 剩余4个元素包含任何阶4 Sylow 子群的全部元素, 所以该 Sylow 子群唯一, 正规.
\end{solution}
\begin{exercise}有限群作用于自身的子群集合. 证明正规化子是子群, 并说明它与中心化子的区别.\end{exercise}
\begin{solution}
正规化子是此作用的稳定子, 故成群. 正规化要求整体送回同一子群, 可以置换其中元素; 中心化要求每个元素不变. 例如 $S_3$ 的 $A_3$ 正规化子是整个 $S_3$, 中心化子只有 $A_3$.
\end{solution}

\originalchapter{环、理想与商环}
\originalsection{两种运算}
\begin{definition}
本书通常讨论有乘法单位元且 $1\ne0$ 的环; 为使商环和同构定理包含全理想情形, 也允许零环作为退化例外, 其 $1=0$. 域和整环始终要求 $1\ne0$. 加法使 $R$ 成交换群, 乘法结合且满足左右分配律. 不额外说“交换”时允许乘法不交换. 环同态保持加法、乘法和单位元. 子环包含同一单位元. 非零元素若有乘法逆元称为单位. 交换环若无非零零因子称为整环; 若每个非零元素可逆称为域.
\end{definition}
整数环、域、多项式环和矩阵环都是例子. $\Z/n\Z$ 中 $\bar a$ 可逆当且仅当 $\gcd(a,n)=1$: 贝祖等式给出逆元, 反之 $ab\equiv1$ 说明任何公因子整除1. 因此 $\Z/n\Z$ 是域当且仅当 $n$ 素数. 非零有限整环总是域: 对 $a\ne0$, 映射 $x\mapsto ax$ 单射, 有限性给出满射, 因而1在其像中.

$0a=(0+0)a=0a+0a$, 加法消去得 $0a=0$, 也有 $a0=0$. 再由 $a+(-a)=0$ 得 $(-a)b=-(ab)$, 最后 $(-a)(-b)=ab$. 这些是分配律的结果, 不能另作乘法公理.
\originalsection{理想与同构定理}
\begin{definition}
环 $R$ 的双边理想 $I$ 是加法子群, 满足 $rI\subseteq I$, $Ir\subseteq I$ 对每个 $r\in R$ 成立. 交换环中只需写 $rI\subseteq I$. 由 $a$ 生成的主理想记 $(a)$; 交换环中 $(a)=\{ra:r\in R\}$.
\end{definition}
理想通常不含1; 若含1则是全环. 环同态的核是理想, 因 $\varphi(ri)=\varphi(r)\varphi(i)=0$, 右乘同理. 对理想 $I$, 在加法陪集上定义 $(a+I)(b+I)=ab+I$: 改变代表元时增量为 $ai+jb+ji\in I$, 故良定义. 此商环 $R/I$ 的单位元为 $1+I$; 真理想保证 $1+I\ne0$. 若 $I=R$, 商为零环, 本书不把它当作域.
\begin{theorem}
对单位保持环同态 $\varphi:R\to S$, 有 $R/\ker\varphi\cong\im\varphi$. 包含 $I$ 的理想与 $R/I$ 的理想由像和原像一一对应. 若 $I\subseteq J$, 则 $(R/I)/(J/I)\cong R/J$.
\end{theorem}
\begin{proof}
映射 $a+\ker\varphi\mapsto\varphi(a)$ 与群情形相同地良定义且双射, 又保持乘法和1. 理想的原像满足加法封闭及两边吸收; 若原理想包含核, 其像在满射的像环中也有吸收性. 原像与像互逆的证明和群对应定理相同. 最后映射 $R/I\to R/J$ 的核为 $J/I$, 用第一结论.
\end{proof}
\originalsection{素理想与极大理想}
\begin{definition}
以下 $R$ 交换. 真理想 $P$ 若 $ab\in P$ 必有 $a\in P$ 或 $b\in P$, 称素理想. 真理想 $M$ 若它和 $R$ 之间无其他理想, 称极大理想.
\end{definition}
\begin{theorem}\label{thm:maximal}
$R/P$ 为整环当且仅当 $P$ 为素理想; $R/M$ 为域当且仅当 $M$ 为极大理想. 因此极大理想必为素理想.
\end{theorem}
\begin{proof}
商环无零因子恰为 $ab\in P$ 迫使一个因子在 $P$ 中. 若 $M$ 极大, 对 $a\notin M$ 有 $M+(a)=R$, 所以存在 $m\in M,r\in R$ 使 $m+ra=1$, 从而 $a+M$ 可逆. 反之若商为域, 对理想 $J\supsetneq M$ 取 $a\in J\setminus M$, 其像可逆, 故 $J/M$ 含1, 即 $J=R$.
\end{proof}
\originalsection{中国剩余定理}
理想 $I,J$ 的和由 $i+j$ 组成, 交为普通集合交, 积 $IJ$ 由有限和 $\sum i_kj_k$ 组成; 逐项分配可证它们都是理想. 总有 $IJ\subseteq I\cap J$, 但一般不相等.
\begin{theorem}[中国剩余]\label{thm:crt}
交换环中若 $I+J=R$, 则 $I\cap J=IJ$, 且 $R/IJ\cong R/I\times R/J$.
\end{theorem}
\begin{proof}
取 $i+j=1$, $i\in I,j\in J$. 对 $x\in I\cap J$, 写 $x=xi+xj\in IJ$, 故交等于积. 映射 $r\mapsto(r+I,r+J)$ 的核是交. 给定 $(a+I,b+J)$, 元素 $aj+bi$ 映到它, 因 $j\equiv1$ 模 $I$ 且 $i\equiv1$ 模 $J$, 所以满射. 同构定理给出结论.
\end{proof}
\begin{example}构造 $\Z/15\Z$ 的非平凡幂等元, 并说明有限环的单位与零因子.\end{example}
\begin{solution}
$\Z/15\Z\cong\F_3\times\F_5$. 幂等元在两个域中各只能0或1, 所以共有4个: $0,1,6,10$. $6^2\equiv6$, $10^2\equiv10$ 模15. 互素于15的元素为单位, 其余非零元素为零因子, 如 $\bar3\bar5=0$.
\end{solution}
\originalsection{带解答练习}
\begin{exercise}证明 $\Z[x]$ 的 $(x)$ 为素理想却不极大, 而 $(2,x)$ 为极大理想.\end{exercise}
\begin{solution}
代入0的满射 $\Z[x]\to\Z$ 的核是 $(x)$, 商为整环而非域. 模2再代入0的满射到 $\F_2$ 核为 $(2,x)$, 商为域.
\end{solution}
\begin{exercise}说明直积环 $R\times S$ 一般不满足任意两环到第三环的余积泛性质 (原附注的坐标注入采用非单位保持约定).\end{exercise}
\begin{solution}
单位保持约定下 $r\mapsto(r,0)$ 不是单位保持映射. 即便允许非单位保持同态, 对 $R=S=T=\Z$ 的两个恒等映射, 想要的加法映射必为 $(a,b)\mapsto a+b$, 但 $(1,0)(0,1)=0$ 的像为0, 两个像的乘积为1. 因而不保持乘法. 原作业8的相关附注须改正; 阿贝尔群和模的有限直和确实与直积相同, 环的情况不能照搬.
\end{solution}

\originalchapter{整除与唯一分解}
\originalsection{分式域与因子}
\begin{theorem}
每个整环 $R$ 嵌入一个分式域 $K$, 其中元素写为 $a/b$, $a,b\in R,b\ne0$, 且任何包含 $R$ 的域都接收唯一延拓 $a/b\mapsto a b^{-1}$.
\end{theorem}
\begin{proof}
在有序对 $(a,b)$ 上定义 $(a,b)\sim(c,d)$ 当且仅当 $ad=bc$. 自反、对称直接; 若 $ad=bc$, $cf=de$, 则 $adf=bde$, 消去非零 $d$ 得 $af=be$, 所以传递. 定义加法 $(a,b)+(c,d)=(ad+bc,bd)$ 和乘法 $(ac,bd)$; 代入等价式可直接核实独立于代表元. 它们满足环公理, 单位为 $(1,1)$, 非零 $(a,b)$ 的逆为 $(b,a)$. 映射 $a\mapsto(a,1)$ 单射. 延拓被 $a/b=a\cdot b^{-1}$ 强制决定, 且交叉乘积等价保证它良定义.
\end{proof}
\begin{definition}
整环中 $a\mid b$ 指 $b=ac$ 对某个 $c\in R$ 成立. 两元素相伴指相差一个单位因子. 非零非单位 $p$ 若 $p=ab$ 迫使 $a$ 或 $b$ 为单位, 称不可约元; 若 $p\mid ab$ 迫使 $p\mid a$ 或 $p\mid b$, 称素元. 每个非零非单位都可分解为不可约元, 且分解除次序与相伴外唯一的整环称唯一分解整环 (UFD).
\end{definition}
素元一定不可约: 若 $p=ab$, 则 $p\mid a$ 或 $p\mid b$; 如 $a=pc$, 消去 $p$ 得 $1=cb$, 所以 $b$ 为单位. 反向一般不成立. 但在 UFD 中不可约元必为素元: 若 $p\mid ab$, 分解 $a,b$ 及商 $ab/p$ 为不可约元, 唯一性迫使 $p$ 与 $a$ 或 $b$ 的某个因子相伴, 故整除其中一个. 零因子输入时结论直接成立. 因而 UFD 中不可约元 $p$ 的商环 $R/(p)$ 为整环.
\originalsection{欧几里得整环与主理想整环}
\begin{definition}
每个理想均主的整环称主理想整环 (PID). 若存在函数 $\delta:R\setminus\{0\}\to\N\cup\{0\}$, 使对 $a,b\in R,b\ne0$ 可写 $a=bq+r$, 且 $r=0$ 或 $\delta(r)<\delta(b)$, 则称欧几里得整环.
\end{definition}
\begin{theorem}\label{thm:pidufd}
欧几里得整环是 PID, PID 是 UFD. PID 中不可约元都是素元, 任意 $a,b$ 的最大公因子 $d$ 可写为 $ra+sb$.
\end{theorem}
\begin{proof}
非零理想中选 $\delta$ 最小的非零元素 $b$. 对 $a\in I$ 除以 $b$ 得余项仍在 $I$, 最小性迫使余项0, 故 $I=(b)$.

在 PID 中 $(a,b)=(d)$, 因 $a,b\in(d)$, $d$ 是公因子; 又 $d=ra+sb$, 任一公因子整除 $d$. 若 $p$ 不可约, 理想 $(p)\subseteq(a)$ 迫使 $a\mid p$, 因而 $(a)=(p)$ 或 $(a)=R$. 所以 $(p)$ 极大, 从定理\ref{thm:maximal}得到 $p$ 素.

为证分解存在, 先证 PID 的理想升链停止. 升链 $I_1\subseteq I_2\subseteq\cdots$ 的并是理想, 可写 $(d)$. 元素 $d$ 属于某一 $I_N$, 从而整个并等于 $I_N$. 若某个非零非单位不能有限分解, 它必可拆为两个非单位, 其中至少一个仍不能有限分解. 反复选此因子 $a_{k+1}$, 得严格升链 $(a_k)\subsetneq(a_{k+1})$, 矛盾. 唯一性则从 $p_1\cdots p_r=q_1\cdots q_s$ 开始, 素性使 $p_1$ 整除某个 $q_j$, 不可约性使它们相伴; 消去后归纳即可.
\end{proof}
\originalsection{高斯整数与分解失败}
高斯整数 $\Z[i]$ 的范数 $N(a+bi)=a^2+b^2$ 为非负整数, 且乘法性由复数绝对值给出. 对 $z,w\in\Z[i],w\ne0$, 将 $z/w$ 的实虚部各舍入至最近整数得 $q$, 则 $|z/w-q|^2\le1/2$, 所以 $r=z-qw$ 满足 $N(r)\le N(w)/2<N(w)$. 故 $\Z[i]$ 是欧几里得整环.
\begin{example}说明 $\Z[\sqrt{-5}]$ 不是 UFD, 尽管每个非零非单位都能分解为不可约元.\end{example}
\begin{solution}
范数 $N(a+b\sqrt{-5})=a^2+5b^2$ 乘法性成立, 值为正整数. 单位仅 $\pm1$. 非单位范数至少2, 所以可用范数归纳完成有限分解. 范数2和3无解, 因而范数4的2、范数9的3、范数6的 $1\pm\sqrt{-5}$ 均不可约: 若分成非单位, 范数因子会要求不存在的2或3. 但
\[
6=2\cdot3=(1+\sqrt{-5})(1-\sqrt{-5}),
\]
两边不可约因子的范数分别为4,9与6,6, 不可能相伴, 唯一性失败.
\end{solution}
\originalsection{带解答练习}
\begin{exercise}证明 UFD 中两个非零元素的 gcd 和 lcm 存在, 且其乘积与 $ab$ 相伴.\end{exercise}
\begin{solution}
将涉及的不同素因子列为 $p_1,\ldots,p_k$, 写 $a=u\prod p_i^{r_i}$, $b=v\prod p_i^{s_i}$. gcd 可取指数 $\min(r_i,s_i)$, lcm 取 $\max(r_i,s_i)$; 任意公因子或公倍数的指数比较给出定义. 因最小与最大之和等于 $r_i+s_i$, 两者乘积与 $ab$ 相伴.
\end{solution}
\begin{exercise}证明 $\Z[x]$ 不是 PID.\end{exercise}
\begin{solution}
若 $(2,x)=(f)$, 则 $f\mid2$, 多项式次数相加迫使 $f$ 为整数常数, 只能相伴于1或2. 若为1, 则1可写 $2a(x)+xb(x)$, 代入0得到1为偶数; 若为2, 则 $2\mid x$, 不成立. 因此该理想不主.
\end{solution}

\originalchapter{多项式与不可约性}
\originalsection{除法与根}
在交换环 $R$ 上, 多项式是有限形式和 $\sum a_ix^i$, 不是仅指其代值函数; 在有限域上不同多项式可能定义同一函数, 如 $x^q-x$ 处处取0. 在整环中非零多项式满足 $\deg(fg)=\deg f+\deg g$, 因最高项系数乘积非零.
\begin{theorem}
若 $F$ 为域, 则任意 $f,g\in F[x]$, $g\ne0$, 存在唯一 $q,r$ 使 $f=qg+r$, $r=0$ 或 $\deg r<\deg g$. 所以 $F[x]$ 是欧几里得整环. 次数 $d$ 的非零多项式在任意包含 $F$ 的域中至多有 $d$ 个不同根.
\end{theorem}
\begin{proof}
若 $\deg f\ge\deg g$, 用最高项之比乘相应 $x$ 幂消去 $f$ 最高项; 次数严格下降, 迭代终止. 两种除法相减时 $(q-q')g=r'-r$, 非零左侧次数至少 $\deg g$, 与右侧冲突, 故唯一. 代入 $a$ 的余项定理给出 $f(a)=0$ 等价于 $(x-a)\mid f$. 每剥去一个根次数降1, 且其他不同根仍是商的根, 归纳得根数界.
\end{proof}
\originalsection{高斯引理}
若 $R$ 为 UFD, 多项式系数的 gcd 称内容, 差一个单位. 内容为单位称本原多项式. 每个非零多项式可写成内容乘本原多项式.
\begin{lemma}[高斯]\label{lem:gauss}
UFD 上两个本原多项式的乘积仍本原. 对分式域 $K$, 本原 $f\in R[x]$ 在 $R[x]$ 不可约当且仅当在 $K[x]$ 不可约. $R[x]$ 本身也是 UFD.
\end{lemma}
\begin{proof}
若某素元 $p$ 整除乘积的全部系数, 则把系数模 $p$ 化到整环 $R/(p)$ 后, 两个非零多项式的乘积为零, 矛盾. 所以乘积本原, 一般内容也满足 $c(fg)\sim c(f)c(g)$.

任意 $K[x]$ 多项式清分母后再除内容可写成 $u g_0$, $u\in K^\times$, $g_0\in R[x]$ 本原. 若本原 $f=gh$ 在 $K[x]$ 非平凡分解, 则 $f=uvg_0h_0$. 将 $uv=a/b$ 写成非零 $a,b\in R$ 的互素分数. 等式 $bf=ag_0h_0$ 的两边内容分别与 $b,a$ 相伴, 因 $f,g_0h_0$ 都本原. 故 $a,b$ 相伴. 互素性迫使两者都是单位, 所以 $uv$ 为 $R$ 的单位. 因而分解可在 $R[x]$ 实现. 反向显然, 而本原多项式的常数非单位因子不可能出现.

最后在 $K[x]$ 唯一分解 $f/c(f)$, 把每个因子改成本原 $R[x]$ 因子, 得存在性. 唯一性由 $K[x]$ 的唯一性和上述内容比较给出, 常数内容再用 $R$ 的唯一性分解. 因而得到 $R[x]$ 的唯一分解.
\end{proof}
\originalsection{约化与艾森斯坦判别}
\begin{proposition}
设 $f\in\Z[x]$ 本原, 模素数 $p$ 后次数不变, 且 $\bar f\in\F_p[x]$ 不可约, 则 $f$ 在 $\Q[x]$ 不可约.
\end{proposition}
\begin{proof}
若可约, 高斯引理给出正次数整数分解 $f=gh$. 最高系数模 $p$ 非零保证两因子的次数都不变, 因而约化给出 $\bar f$ 的非平凡分解, 矛盾.
\end{proof}
\begin{theorem}[艾森斯坦]\label{thm:eisenstein}
设 $f=a_nx^n+\cdots+a_0\in\Z[x]$, $n\ge1$, 存在素数 $p$ 使 $p\nmid a_n$, $p\mid a_i$ 对 $i<n$ 成立, 且 $p^2\nmid a_0$. 则 $f$ 在 $\Q[x]$ 不可约.
\end{theorem}
\begin{proof}
先除内容, 条件仍成立. 若 $f=gh$ 为两个正次数整数多项式乘积, 约化后乘积仅为单项式 $\bar a_nx^n$. 在 $\F_p[x]$ 的唯一分解中因子都只能是非零常数乘 $x$ 的幂. 两因子的次数约化不变, 故都含正次 $x$, 两个常数项都被 $p$ 整除. 于是 $p^2\mid a_0$, 矛盾.
\end{proof}
\begin{example}证明 $1+x+\cdots+x^{p-1}$ 对素数 $p$ 在 $\Q[x]$ 不可约.\end{example}
\begin{solution}
代入 $x=y+1$, 得 $((y+1)^p-1)/y=\sum_{k=1}^p\binom pk y^{k-1}$. 除最高系数1外均被 $p$ 整除, 常数 $p$ 不被 $p^2$ 整除. 艾森斯坦判别给出代入后不可约, 而代入 $y=x-1$ 是逆变换, 故原式不可约.
\end{solution}
\originalsection{带解答练习}
\begin{exercise}证明 $x^3-2$ 在 $\Q[x]$ 不可约, 而 $x^4+4$ 可约.\end{exercise}
\begin{solution}
前者对素数2用艾森斯坦. 后者为 $(x^2-2x+2)(x^2+2x+2)$, 由 $(x^2+2)^2-(2x)^2$ 展开可核.
\end{solution}
\begin{exercise}求 $\F_2$ 上所有首一不可约二次、三次多项式.\end{exercise}
\begin{solution}
二次或三次可约当且仅当有一次因子, 即有域内根. 首一候选常数必须1; 再代入1排除. 二次仅 $x^2+x+1$; 三次仅 $x^3+x+1$ 和 $x^3+x^2+1$. 它们在0,1都不取0.
\end{solution}

\originalchapter{模与线性结构}
\originalsection{从向量空间到模}
\begin{definition}
设 $R$ 是交换环. $R$ 模 $M$ 是交换加法群, 带标量乘法 $R\times M\to M$, 满足 $r(m+n)=rm+rn$, $(r+s)m=rm+sm$, $(rs)m=r(sm)$, $1m=m$. 保持加法与标量乘法的映射称模同态. 子模是包含0且对加法、负元及标量乘法封闭的子集.
\end{definition}
域上的模就是向量空间. $\Z$ 模就是交换群: $nm$ 由重复加法及负号强制确定. 环的理想是 $R$ 自身作为模的子模. 任意矩阵 $A\in M_{s\times t}(R)$ 给出同态 $R^t\to R^s$, 核描述线性关系, 商 $R^s/\im A$ 描述按这些关系识别生成元的模.

模的核、像、商与同构定理都按加法群构造, 还需核验标量相容. 若 $m-m'\in N$, 则 $r(m-m')\in N$, 所以 $r(m+N)=rm+N$ 良定义. 对同态 $f$, 映射 $M/\ker f\to\im f$, $m+\ker f\mapsto f(m)$ 双射且线性. 包含 $N$ 的子模和 $M/N$ 的子模由像、原像对应.
\originalsection{自由模与有限生成}
两个模的直和 $A\oplus C$ 是笛卡尔积上的逐坐标加法和标量乘法. 在同一个模中的内部直和指两子模的和为所讨论的模, 且交为0; 每个元素因而有唯一两分量表示.

\begin{definition}
有限列 $m_1,\ldots,m_s$ 生成 $M$ 指每个元素可写 $\sum r_im_i$. 若这种写法唯一, 则该列称基, $M$ 称自由模, 并同构于 $R^s$. 元素 $m$ 的湮灭理想为 $\ann(m)=\{r:rm=0\}$; 由某个非零标量杀掉的元素称挠元, 此处 $R$ 为整环.
\end{definition}
域上每个有限生成模有基; 一般环上不能这样推断. $\Z/2\Z$ 由一个元素生成却不是自由 $\Z$ 模, 因生成元被2杀掉. 对整环, 挠元构成子模: 若 $am=0,bn=0$ 且 $a,b\ne0$, 则 $ab(m+n)=0$; $ab\ne0$ 是这里整环条件的用途.
\begin{theorem}\label{thm:subfree}
PID $R$ 上, $R^n$ 的任何子模自由且可由至多 $n$ 个元素生成.
\end{theorem}
\begin{proof}
按 $n$ 归纳. 对 $n=1$, 子模是理想 $(d)$; 若非零, $r\mapsto rd$ 因整环消去律给出 $R\cong(d)$; 零模以空基表示.

对一般 $n$, 投影到第一坐标的像是理想 $(d)$. 若 $d=0$, 子模位于 $0\times R^{n-1}$, 用归纳. 若 $d\ne0$, 选 $v$ 的第一坐标为 $d$. 任意 $w$ 的第一坐标是 $ad$, 故 $w-av$ 属投影核 $K$. 因 $av$ 第一坐标 $ad$, $Rv\cap K=0$. 于是子模为 $Rv\oplus K$, 核由至多 $n-1$ 个元素自由生成, 加上 $v$ 即为所求基.
\end{proof}
以下 $R$ 为 PID. 有限生成模 $M$ 是某个 $R^s$ 的商. 取生成元给满射 $R^s\to M$, 其核按定理\ref{thm:subfree}有限生成且自由, 选其基得到矩阵 $A$ 与表示 $M\cong R^s/\im A$. 下一章会把矩阵关系化成一维关系, 这正是模理论引入的目的.
\originalsection{正合与分裂}

\begin{definition}
列 $0\to A\xrightarrow{i}B\xrightarrow{p}C\to0$ 称短正合列, 若 $i$ 单射、$p$ 满射且 $\ker p=\im i$. 若有同态 $s:C\to B$ 使 $ps$ 为恒等, 称分裂.
\end{definition}
\begin{proposition}
分裂短正合列满足 $B\cong A\oplus C$. 若 $C$ 自由则一定分裂.
\end{proposition}
\begin{proof}
映射 $(a,c)\mapsto i(a)+s(c)$ 是同态. 对 $b$, 令 $c=p(b)$, 则 $b-s(c)\in\ker p$, 给出满射; 若 $i(a)+s(c)=0$, 应用 $p$ 得 $c=0$, 再由 $i$ 单射得 $a=0$. 若 $C$ 自由, 给每个基向量任选 $B$ 中原像, 唯一线性延拓得到 $s$.
\end{proof}
\originalsection{带解答练习}
\begin{exercise}短正合列 $0\to\Z\xrightarrow{\times2}\Z\to\Z/2\Z\to0$ 是否分裂?\end{exercise}
\begin{solution}
若有截面 $s$, 因 $2\bar1=0$, 必有 $2s(\bar1)=0$ 在 $\Z$ 中, 故 $s(\bar1)=0$, 不能映回 $\bar1$. 不分裂. 这说明不能把所有商模都写成原模的直和因子.
\end{solution}
\begin{exercise}对 $R=\Z$ 和 $M=\Z/12\Z$, 求 $\ann(\bar4)$ 及其生成子模.\end{exercise}
\begin{solution}
$n\bar4=0$ 当且仅当 $3\mid n$, 所以湮灭理想是 $3\Z$. 子模 $\{\bar0,\bar4,\bar8\}$ 同构于 $\Z/3\Z$.
\end{solution}

\originalchapter{Smith 标准形与模分解}
\originalsection{关系矩阵的化简}
给关系矩阵换行对应改变自由模 $R^s$ 的基, 换列对应改变关系的生成列, 两者都不改变商模的同构类型. 允许的初等变换包括交换行列、某一行列乘单位、给一行列加另一行列的倍数, 它们均可逆.
\begin{theorem}[Smith 标准形]\label{thm:smith}
PID $R$ 上的矩阵 $A$ 可经可逆行列变换化为对角矩阵, 非零对角元 $d_1,\ldots,d_r$ 满足 $d_1\mid d_2\mid\cdots\mid d_r$.
\end{theorem}
\begin{proof}
零矩阵已经是所需对角形. 以下非零情形先说明两元素的消元, 其中 $(a,b)\ne(0,0)$. 对 $a,b$ 取 gcd $d=ua+vb$, 写 $a=da'$, $b=db'$, 则 $ua'+vb'=1$. 矩阵
\[
\begin{pmatrix}u&v\\-b'&a'\end{pmatrix}
\]
行列式为1, 将列向量 $(a,b)^T$ 变成 $(d,0)^T$. 同理可作列变换.

把一个非零元素移到左上角. 若首列某元素不能被左上角元素整除, 用上述 gcd 变换把左上角改成严格较大的主理想生成元; 首行也同样处理. 这种严格增大不可能无限继续, 因 PID 的理想升链停止. 因而终可使左上角整除首行首列所有项, 再用减去倍数清零, 得 $\operatorname{diag}(d,B)$.

若 $B$ 中仍有项 $b$ 不被 $d$ 整除, 将该项所在行加到首行, 首行在该列出现 $b$, 再用 gcd 列消元, 左上角主理想又严格增大. 升链停止保证有限次后 $d$ 整除 $B$ 中每项. 对 $B$ 归纳. 在其内部的行列变换仍使每项被 $d$ 整除, 故其第一个非零对角元也被 $d$ 整除. 归纳给出整除链.
\end{proof}
\originalsection{有限生成模结构定理}
\begin{theorem}\label{thm:module}
PID 上任何有限生成模同构于
\[
R^t\oplus R/(d_1)\oplus\cdots\oplus R/(d_r),\qquad d_1\mid\cdots\mid d_r,
\]
其中 $d_i$ 为非零非单位. 自由秩 $t$ 和各 $d_i$ 的相伴类唯一. 等价地, 挠部分可按素元幂分解为初等因子.
\end{theorem}
\begin{proof}
由上一章, $M$ 有有限关系矩阵. Smith 变换给出直和商: 非零对角元 $d$ 给 $R/(d)$, 零行给自由因子, 单位对角元给零模, 可删去.

对唯一性, 先把模延标量到分式域, 每个非零 $d_i$ 在其中可逆, 所以维数为 $t$, 自由秩唯一. 对素元 $p$, 将挠模中的 $p$ 主部分记为 $T_p=\{m:p^k m=0\text{ 对某 }k\}$. 对互素因子的中国剩余定理说明挠模是这些主部分的直和. 在 $T_p$ 上, $p^{k-1}T_p/p^kT_p$ 是域 $R/(p)$ 上的向量空间. 一个因子 $R/(p^a)$ 在该维数中贡献1当且仅当 $a\ge k$, 记该维数为 $b_k$, 则次数恰为 $k$ 的初等因子有 $b_k-b_{k+1}$ 个, 所以这些维数恢复每个素元幂指数的多重集. 最后将每种素元的指数按非降序排列令 $r$ 为各素元初等因子数的最大值 (挠模为0时取 $r=0$), 并在左侧补零至长度 $r$, 原整除链中首个非单位因子必含某素元, 该素元因整除链而出现在全部因子中, 因而原非单位因子数就是此最大值 $r$. 逐位置相乘, 恢复整除链的 $d_i$; 这也证明初等因子与不变因子之间转换唯一.
\end{proof}
这里延标量可具体理解为形式分数 $m/s$, $s\ne0$, 以 $m/s=m'/s'$ 当某非零 $u$ 杀掉 $s'm-sm'$ 来识别; 具体为 $m/s+m'/s'=(s'm+sm')/(ss')$, $(a/b)(m/s)=am/(bs)$. 等价式传递性通过相乘非零分母及两个原见证因子得到新见证; 改变代表元后, 加法或标量乘法两边的交叉乘积之差是原零差的标量线性组合, 再乘原见证因子的乘积即可杀掉, 故运算良定义. 向量空间各公理由分配律和分式运算直接继承. 模同构 $f$ 延拓为 $m/s\mapsto f(m)/s$, 其逆由 $f^{-1}$ 延拓, 因而保留延标量维数. 对直和分量逐坐标取分数即可得到相应直和. 其核恰为挠元, 对自由因子得到 $K^t$, 因而上述自由秩论证不需要预设张量积理论.
\begin{example}分解由矩阵 $\begin{pmatrix}4&6\\2&8\end{pmatrix}$ 给出的 $\Z$ 模.\end{example}
\begin{solution}
所有元素的 gcd 是2, 行列式为20. Smith 形第一个因子为2, 第二个为10: 可先交换行, 将第二行减第一行2倍得 $(0,-10)$, 首行是 $(2,8)$, 再将第二列减第一列4倍清掉8, 最后改变符号. 商为 $\Z/2\Z\oplus\Z/10\Z$, 阶20.
\end{solution}
\originalsection{线性算子的代数解释}
给域 $F$ 上有限维空间 $V$ 和线性算子 $T$, 定义 $f(x)v=f(T)v$, 使 $V$ 成 $F[x]$ 模. 它全为挠模: 任意 $v$, 向量 $v,Tv,\ldots,T^nv$ 线性相关, 所以某非零多项式杀掉 $v$. 结构定理给出循环因子 $F[x]/(f_i)$, $f_i\mid f_{i+1}$. 在基 $1,x,\ldots,x^{d-1}$ 中, 乘 $x$ 的矩阵是 $f_i$ 的伴随矩阵, 得到有理标准形. 若多项式分裂, 初等因子 $F[x]/((x-\lambda)^a)$ 在基 $1,(x-\lambda),\ldots$ 中给出一个 Jordan 块. 此处是模分解的应用, 不展开数值计算或稳定性问题.
\originalsection{带解答练习}
\begin{exercise}列出阶72的所有交换群同构类型.\end{exercise}
\begin{solution}
$72=2^3 3^2$. 2主部分为 $C_8,C_4\times C_2,C_2^3$; 3主部分为 $C_9,C_3^2$. 任取两类直积, 共6种. 唯一性由初等因子保证, 所以不会重复.
\end{solution}
\begin{exercise}若 $T^2=0$ 且 $\dim V=5$, $\operatorname{rank}T=2$, 求 Jordan 块.\end{exercise}
\begin{solution}
仅能有大小1或2的零特征值块, 每个大小2块秩为1. 因而有两个2块, 剩余一个1块. 无须假设底域代数闭, 因最小多项式已经分裂为 $x^2$ 的因子.
\end{solution}

\originalchapter{域扩张与代数元}
\originalsection{次数与最小多项式}
\begin{definition}
若域 $K$ 包含域 $F$, 称 $K/F$ 为域扩张, 次数 $[K:F]$ 是 $K$ 作为 $F$ 向量空间的维数. 若 $\alpha\in K$ 是某非零 $F$ 多项式的根, 称其对 $F$ 代数; 否则称超越. 最小包含 $F$ 和 $\alpha$ 的子域记 $F(\alpha)$.
\end{definition}
\begin{theorem}\label{thm:minpoly}
代数元 $\alpha$ 有唯一首一最小多项式 $m_\alpha\in F[x]$, 它不可约, 且
$F(\alpha)=F[\alpha]\cong F[x]/(m_\alpha)$, 次数为 $\deg m_\alpha$.
\end{theorem}
\begin{proof}
代值同态 $F[x]\to K$ 的核非零, 在 PID 中有唯一首一生成元 $m$. 若 $m=gh$ 正次数分解, 则 $g(\alpha)h(\alpha)=0$, 域无零因子迫使低次因子也在核中, 矛盾. 不可约元在 $F[x]$ 中生成极大理想, 故像 $F[\alpha]$ 已是域, 就等于 $F(\alpha)$. 除法给每个余数类唯一代表次数小于 $\deg m$, 所以 $1,\alpha,\ldots,\alpha^{d-1}$ 为基.
\end{proof}
\begin{theorem}[塔公式]\label{thm:tower}
若 $F\subseteq E\subseteq K$ 且两段有限, 则 $[K:F]=[K:E][E:F]$.
\end{theorem}
\begin{proof}
取 $E/F$ 基 $u_i$ 和 $K/E$ 基 $v_j$. 先按 $v_j$ 再按 $u_i$ 展开说明 $u_iv_j$ 张成 $K$; 若 $\sum_{i,j}a_{ij}u_iv_j=0$, 按 $v_j$ 独立得每个 $\sum_i a_{ij}u_i=0$, 再按 $u_i$ 独立得全部系数0.
\end{proof}
有限扩张中每个元素代数, 因幂列线性相关. 有限个代数元逐次添加形成有限扩张, 故代数元的和、积、差及非零逆元仍代数. 塔公式也说明: 对中间域 $E$, $[E:F]$ 整除 $[K:F]$.
\originalsection{根的添加与分裂域}
任何不可约 $f\in F[x]$ 可以通过域 $F[x]/(f)$ 添加一个根 $\bar x$. 对一般非零多项式先添加一个不可约因子的根, 除去一次因子, 再归纳; 经过有限次扩张, 全部根都存在.
\begin{definition}
多项式 $f$ 在 $F$ 上的分裂域是由其全部根生成的域, 且 $f$ 在其中完全分裂为一次因子.
\end{definition}
\begin{lemma}[嵌入延拓]\label{lem:extend}
若 $\sigma:F\to F'$ 为域同构, $\alpha$ 对 $F$ 的最小多项式为 $m$, 则将 $\alpha$ 送至 $\sigma(m)$ 的任意根 $\beta$, 可唯一延拓为 $F(\alpha)\cong F'(\beta)$. 因而分裂域除保持底域的同构外唯一.
\end{lemma}
\begin{proof}
同构将系数逐个变换, $m$ 不可约故 $\sigma(m)$ 不可约. 两边都同构于相应多项式商域, 对应 $\sum a_i\alpha^i\mapsto\sum\sigma(a_i)\beta^i$, 因模最小多项式识别而良定义. 对分裂域, 逐个延拓根的像; 目标域含有对应多项式所有根, 每步可以继续. 最终根全被送入目标分裂域, 其像包含全部根, 故满射.
\end{proof}
\originalsection{可分性}
形式导数定义为 $f'=\sum i a_ix^{i-1}$, 满足乘法法则. 在分裂域中一个根重数至少2当且仅当同时是 $f$ 和 $f'$ 的根; 由 $(x-a)^rg$ 直接求导可核. 所以 $f$ 无重根当且仅当 $\gcd(f,f')=1$. 代数元称可分指最小多项式无重根; 有限扩张可分指每个元素可分.

特征0的不可约多项式导数非零且次数降低, gcd 必为1. 特征 $p$ 时导数为0恰为所有指数被 $p$ 整除, 此时可能不可分. 例如在 $\F_p(t)$ 上 $x^p-t$ 不可约: 若 $t$ 是 $p$ 次幂, 分子分母的 $t$ 指数比较将矛盾; 更完整地, 在 $\F_p[t][x]$ 对素元 $t$ 用艾森斯坦证明不可约. 这里采用 UFD 版本: 若素元整除全部低次系数而其平方不整除常数项, 最高系数不被它整除, 则分式域上不可约. 第7章证明逐字适用: 高斯引理清分母, 再在整环 $R/(p)$ 中使用 $x$ 是素元且多项式次数相加, 推得两因子常数项均被 $p$ 整除, 造成平方整除的矛盾. 因此不依赖未经证明的推广. 它只有一个不同根, 重数 $p$.
\begin{example}求 $\Q(\sqrt2,\sqrt3)$ 的次数和一组基.\end{example}
\begin{solution}
$\sqrt3\notin\Q(\sqrt2)$: 若 $\sqrt3=a+b\sqrt2$, $a,b\in\Q$, 平方给 $3=a^2+2b^2+2ab\sqrt2$. 因 $\sqrt2$ 无理, $ab=0$; $b=0$ 或 $a=0$ 都迫使有理数平方为3或 $3/2$, 用素因子奇偶性可排除. 所以两次扩张各次数2, 总次数4, 基为 $1,\sqrt2,\sqrt3,\sqrt6$.
\end{solution}
\originalsection{带解答练习}
\begin{exercise}在 $\Q(\sqrt2)$ 中求 $(1+\sqrt2)^{-1}$, 并求 $\sqrt2+\sqrt3$ 的最小多项式.\end{exercise}
\begin{solution}
逆元为 $\sqrt2-1$. 令 $\alpha=\sqrt2+\sqrt3$, 则 $\alpha^2=5+2\sqrt6$, 所以 $\alpha^4-10\alpha^2+1=0$. 又 $\alpha^{-1}=\sqrt3-\sqrt2$, 故两者和差恢复 $\sqrt2,\sqrt3$, $\Q(\alpha)$ 即次数4的双二次域, 最小多项式因此为该首一四次式.
\end{solution}

\originalchapter{有限域}
\originalsection{特征与 Frobenius}
\begin{definition}
域 $F$ 的特征是 $n\cdot1=0$ 的最小正整数, 不存在时记为0. 特征若正则必为素数 $p$: 若 $n=ab$ 合成, $(a\cdot1)(b\cdot1)=0$ 与最小性冲突. $F$ 的最小子域为特征0时的 $\Q$, 或特征 $p$ 时的 $\F_p$.
\end{definition}
有限域是其素域上的有限维向量空间, 因而元素数必为 $p^n$. 在特征 $p$ 的交换环中, 二项式中间系数被 $p$ 整除, 所以 $(a+b)^p=a^p+b^p$. 映射 $\mathrm{Fr}:a\mapsto a^p$ 保持加法、乘法和1, 在域上单射; 有限性使其满射, 因而为自同构.
\begin{proposition}
有限域上每个不可约多项式可分.
\end{proposition}
\begin{proof}
若不可约 $f$ 不可分, $f'$ 必为0, 故 $f=\sum a_ix^{pi}$. Frobenius 满射使每个系数 $a_i=b_i^p$, 从而 $f=(\sum b_ix^i)^p$, 不可约性矛盾.
\end{proof}
\originalsection{存在性与唯一性}
\begin{theorem}\label{thm:finitefield}
对每个素数幂 $q=p^n$, 存在唯一 (除同构外) 阶 $q$ 的域, 记为 $\F_q$. 它是 $x^q-x$ 在 $\F_p$ 上的分裂域.
\end{theorem}
\begin{proof}
取 $x^q-x$ 的分裂域 $L$. 导数为 $-1$, 故有恰好 $q$ 个不同根. 根集合 $S$ 对加减乘法封闭, 因 $(a\pm b)^q=a^q\pm b^q$, $(ab)^q=a^qb^q$; 对非零 $a$, $(a^{-1})^q=(a^q)^{-1}=a^{-1}$. 且0,1在其中, 所以 $S$ 是域. 它包含全部根, 分裂域的最小性给 $L=S$, 故阶 $q$.

反之任意阶 $q$ 域的非零乘法群阶 $q-1$, 拉格朗日定理给 $a^{q-1}=1$, 所有元素满足 $a^q=a$. 因此它就是 $x^q-x$ 的分裂域, 唯一性来自分裂域唯一性.
\end{proof}
\originalsection{乘法群与子域}
\begin{theorem}\label{thm:finitecyclic}
域的有限乘法子群 $H$ 是循环群.
\end{theorem}
\begin{proof}
$H$ 交换. 令各元素阶的最小公倍数为 $m$. 对每个素数幂 $p^a\mid m$ 的最高幂, 有某个元素阶含该幂; 取其适当幂得到阶恰为 $p^a$ 的元素. 这些元素相乘, 因交换且阶互素, 乘积的阶为各阶之积 $m$: 若乘积某次幂为1, 一个因子幂等于其余因子幂的逆, 两边阶互素, 故每一因子幂均为1. 因而 $H$ 有阶 $m$ 元素. 同时每个 $h$ 满足 $h^m=1$, 多项式根数界给 $|H|\le m$. 而该元素已经给出 $m$ 个不同幂, 所以 $|H|=m$, 群循环.
\end{proof}
$\F_{p^n}$ 上 Frobenius 的阶恰为 $n$: $n$ 次迭代为恒等, 若 $0<d<n$ 次迭代恒等, 则 $x^{p^d}-x$ 会有 $p^n$ 个根, 超过次数. 对 $d\mid n$, 多项式 $x^{p^d}-x$ 的根全在 $\F_{p^n}$ 中: 非零根满足阶整除 $p^d-1$, 而 $p^d-1\mid p^n-1$, 循环乘法群中恰有 $p^d-1$ 个这样的元素. 加0得唯一 $p^d$ 元子域. 任何子域的次数 $d$ 整除 $n$, 由塔公式, 且其元素必为该多项式的根, 故没有其他子域.
\begin{example}在 $\F_2[t]/(t^3+t+1)$ 中证明 $\alpha=\bar t$ 生成乘法群.\end{example}
\begin{solution}
三次式无 $\F_2$ 根所以不可约, 商域阶8. $\alpha\ne0,1$, 而乘法群阶7是素数, 因此 $\alpha$ 阶7, 生成整个乘法群. 可用 $\alpha^3=\alpha+1$ 逐次约化检查全部七个非零元素.
\end{solution}
\originalsection{带解答练习}
\begin{exercise}列出 $\F_{p^{12}}$ 的所有子域及其包含关系.\end{exercise}
\begin{solution}
次数为 $12$ 的正因子 $1,2,3,4,6,12$, 分别给 $\F_p,\F_{p^2},\F_{p^3},\F_{p^4},\F_{p^6},\F_{p^{12}}$. 包含当且仅当次数整除: 必要性由塔公式, 充分性由相应根集合包含.
\end{solution}
\begin{exercise}证明 $\F_p(t)$ 的 Frobenius 不是满射.\end{exercise}
\begin{solution}
若 $(f(t)/g(t))^p=t$, 约成互素分子分母, 则 $f^p=tg^p$. 素因子 $t$ 在左侧的指数是 $p$ 的倍数, 右侧是 $1$ 加 $p$ 的倍数, 矛盾. 因而 $t$ 没有原像.
\end{solution}

\originalchapter{自同构与固定域}
\originalsection{嵌入数与正规扩张}
$F$ 嵌入指固定 $F$ 的单射域同态. 在下面的有限扩张计数中, 固定目标域 $\Omega$, 为所有选定生成元的 $F$ 最小多项式乘积的分裂域; 它容纳所有需要的共轭根.
底域 $F$ 保持不动的 $K$ 自同构组成群 $\Aut_F(K)$, 记作 $\Gal(K/F)$ 时并不自动断言扩张是 Galois. 任一自同构将代数元送至同一最小多项式的一个根 (允许为自身), 因为逐系数代入保持0.
\begin{proposition}\label{prop:embeddings}
有限扩张 $K/F$ 向包含相关分裂域的域内的 $F$ 嵌入数至多为 $[K:F]$. 若 $K$ 由可分元生成, 则恰为 $[K:F]$, 并且扩张的每个元素可分.
\end{proposition}
\begin{proof}
有限扩张可由有限多个代数元逐次生成. 每一步延拓数至多为该元在当前域上最小多项式的次数, 引理\ref{lem:extend}和塔公式给出上界. 若生成元对 $F$ 可分, 它在当前域上的最小多项式整除其 $F$ 最小多项式, 仍无重根, 运输后的当前最小多项式整除相应生成元的 $F$ 最小多项式, 所以仍在 $\Omega$ 完全分裂; 每个不同根给一个延拓, 所以每一步延拓数都等于次数, 总数相乘恰为 $[K:F]$.

若 $\beta\in K$, 所有嵌入先限制到 $F(\beta)$, 再延拓到 $K$. 若 $\beta$ 最小多项式的不同根少于其次数, 计数上界将严格小于 $[F(\beta):F][K:F(\beta)]$, 与已经得到的总数矛盾. 故每个 $\beta$ 可分.
\end{proof}
\begin{definition}
有限扩张 $K/F$ 正规指任何在 $K$ 中有根的 $F$ 不可约多项式在 $K$ 中完全分裂. 既正规又可分称为有限 Galois 扩张.
\end{definition}
有限扩张正规当且仅当它是有限多个 $F$ 多项式的分裂域. 一方向由有限生成元的最小多项式全部分裂, 它们的根生成的域就是 $K$. 另一方向可用嵌入延拓: 若 $K$ 是分裂域, 为核对另外一个不可约多项式时, 将目标取为它与原生成多项式的共同分裂域. 任一嵌入把生成根排列到同一根集合, 故像为 $K$; 对任意有根的不可约式, 先将该根送至任意其他根, 再逐次延拓到 $K$, 所以其他根也在 $K$ 中. 这同时证明正规等价于所有 $F$ 嵌入均保持 $K$ 本身. 因而有限 Galois 扩张满足 $|\Gal(K/F)|=[K:F]$.
\originalsection{Artin 固定域定理}
\begin{lemma}[不同域同态的独立性]\label{lem:independent}
从域 $L$ 到域 $L'$ 的不同同态 $\sigma_1,\ldots,\sigma_n$ 作为函数在 $L'$ 上线性无关.
\end{lemma}
\begin{proof}
若存在非零线性关系, 选非零项数最小者, 归一化最后系数为1. 只有一项不可能, 因在1处取值非零. 取 $a$ 使 $\sigma_1(a)\ne\sigma_n(a)$. 把关系在 $ax$ 处的式减去 $\sigma_n(a)$ 乘在 $x$ 处的式, 最后一项消失, 第一项不消失, 得到项数更少的非零关系, 矛盾.
\end{proof}
\begin{theorem}[Artin]\label{thm:artin}
设 $H$ 为域 $L$ 的有限自同构群, $E=L^H=\{x:\sigma(x)=x\text{ 对所有 }\sigma\in H\}$. 则 $L/E$ 有限 Galois, $[L:E]=|H|$, 且 $\Gal(L/E)=H$.
\end{theorem}
\begin{proof}
自同构保持加减乘法和非零逆元, 所以共同固定集合 $E$ 确为子域. 先证次数上界. 设 $|H|=n$, 将列向量 $v(x)=(\sigma(x))_{\sigma\in H}\in L^n$ 的所有值取张成空间, 选基列 $v(x_1),\ldots,v(x_r)$, $r\le n$. 对任意 $x$, 有唯一系数 $c_i\in L$ 使 $v(x)=\sum c_iv(x_i)$. 对等式全部作用 $\tau\in H$, 其行指标从 $\sigma$ 变成 $\tau\sigma$, 只是排列, 所以 $v(x)=\sum\tau(c_i)v(x_i)$. 唯一性给 $\tau(c_i)=c_i$, 即 $c_i\in E$. 在恒等同态行取值得 $x=\sum c_ix_i$, 故 $[L:E]\le r\le n$.

反之令 $d=[L:E]$, 每个 $E$ 线性函数 $L\to L$ 由一组 $E$ 基的 $d$ 个取值决定, 这些函数作为 $L$ 向量空间维数 $d$. $H$ 中的 $n$ 个函数由引理\ref{lem:independent}独立, 所以 $n\le d$, 得 $d=n$.

对任意 $x\in L$, 取其在 $H$ 下不同像的乘积多项式 $\prod_y(t-y)$. $H$ 排列这些根, 所以系数在 $E$. 该式无重根且全在 $L$ 分裂, 最小多项式整除它, 从而每个元素可分, 并且其最小多项式全分裂, 得正规. 最后 $H\subseteq\Gal(L/E)$, 后者大小至多 $[L:E]=n$, 所以相等.
\end{proof}
\originalsection{本原元}
\begin{theorem}
有限可分扩张由一个元素生成.
\end{theorem}
\begin{proof}
有限底域时扩域也有限, 非零乘法群循环, 其生成元同时生成整个域. 无限底域时, 先处理 $F(\alpha,\beta)$. 设其 $F$ 嵌入为 $\sigma_1,\ldots,\sigma_d$. 选 $c\in F$ 使所有 $\sigma_i(\alpha+c\beta)$ 两两不同: 对每对不同嵌入, 若 $\beta$ 像不同至多禁掉一个 $c$, 若 $\beta$ 像相同则 $\alpha$ 像不同, 不禁任何 $c$. 仅有限个禁值, 无限域中可选. 元素 $\gamma=\alpha+c\beta$ 有至少 $d$ 个不同共轭, 最小多项式次数至少 $d$, 又不超过 $[F(\alpha,\beta):F]=d$, 所以 $F(\gamma)$ 就是整个域. 对有限多个生成元归纳.
\end{proof}
\originalsection{带解答练习}
\begin{exercise}为什么 $\Q(\sqrt[3]2)/\Q$ 不是 Galois?\end{exercise}
\begin{solution}
$x^3-2$ 不可约且特征0可分, 但另两根非实, 不在这个实域中, 所以不正规. 自同构必须把实根送至同域内的根, 只有自身, 因而自同构群平凡, 大小小于次数3.
\end{solution}

\originalchapter{Galois 对应与计算}
\originalsection{基本定理}
\begin{theorem}[有限 Galois 对应]\label{thm:galois}
设 $L/F$ 为有限 Galois 扩张, $G=\Gal(L/F)$. 中间域 $E$ 与子群 $H\le G$ 由
$E\mapsto\Gal(L/E)$ 和 $H\mapsto L^H$ 互逆对应, 且反转包含关系. 有 $[L:E]=|H|$, $[E:F]=[G:H]$. 中间域 $E/F$ 正规当且仅当 $H$ 在 $G$ 中正规, 此时 $\Gal(E/F)\cong G/H$.
\end{theorem}
\begin{proof}
对任意中间域, $L/E$ 仍是可分多项式的分裂域 (把原 $F$ 多项式看作 $E$ 多项式), 所以为 Galois, 自同构数等于 $[L:E]$. 记 $H=\Gal(L/E)$, Artin 定理给 $[L:L^H]=|H|=[L:E]$, 而 $E\subseteq L^H$, 塔公式迫使相等. 对任意 $H$, Artin 定理直接给 $\Gal(L/L^H)=H$. 包含反转来自“固定更多自同构则元素更少”, 次数公式由塔公式.

对 $g\in G$, 直接检查 $g(L^H)=L^{gHg^{-1}}$. 若 $H$ 正规, 则所有 $g$ 保持 $E=L^H$; 任意 $F$ 嵌入 $E$ 可逐根延拓为 $L$ 的自同构, 因 $L/F$ 正规, 所以所有嵌入均保持 $E$, $E/F$ 正规. 反之 $E/F$ 正规使所有 $g$ 保持 $E$, 上述固定域等式与对应唯一性给 $gHg^{-1}=H$. 可分性继承, 故 $E/F$ Galois. 限制同态 $G\to\Gal(E/F)$ 的核为 $H$, 满射由嵌入延拓, 第一同构定理给商群结论.
\end{proof}
\originalsection{双二次扩张}
\begin{example}求 $L=\Q(\sqrt2,\sqrt3)$ 的所有中间域.\end{example}
\begin{solution}
它是 $(x^2-2)(x^2-3)$ 的分裂域, 次数4. 独立改变两个平方根符号得到四个自同构, 群为 $C_2\times C_2$. 三个阶2子群分别固定 $\sqrt2,\sqrt3,\sqrt6$, 固定域为对应二次域, 因次数都是2. 加上 $\Q$ 和 $L$ 共五个中间域, 没有其他.
\begin{center}
\begin{tikzpicture}[every node/.style={font=\normalsize},>=Stealth]
\node (t) at (0,2.2) {$\Q(\sqrt2,\sqrt3)$};
\node (a) at (-3,1.1) {$\Q(\sqrt2)$};
\node (b) at (0,1.1) {$\Q(\sqrt3)$};
\node (c) at (3,1.1) {$\Q(\sqrt6)$};
\node (d) at (0,0) {$\Q$};
\draw (t)--(a)--(d) (t)--(b)--(d) (t)--(c)--(d);
\end{tikzpicture}
\end{center}
图中的线表示域包含关系, 自同构子群的包含次序反向.
\end{solution}
\originalsection{三次多项式的分裂域}
\begin{example}求 $x^3-2$ 的分裂域、Galois 群及中间域.\end{example}
\begin{solution}
令 $\alpha=\sqrt[3]2$, $\omega$ 为非实三次单位根. 分裂域 $L=\Q(\alpha,\omega)$, $\Q(\alpha)$ 次数3且实, $\omega$ 不在其中而满足二次式 $x^2+x+1$, 故总次数6. 自同构置换三个根, 给出 $G\hookrightarrow S_3$; Galois 性使 $|G|=6$, 所以 $G=S_3$.

$S_3$ 的一个阶3子群给唯一二次域 $\Q(\omega)$, 它正规. 三个阶2子群给三个三次域 $\Q(\alpha),\Q(\omega\alpha),\Q(\omega^2\alpha)$; 各根稳定子为阶2, 固定域次数3, 所以恰是这些域, 彼此不同. 再加端点 $\Q,L$ 共六个中间域. 三个阶2子群不正规, 故三个三次域在 $\Q$ 上不正规.
\end{solution}
\originalsection{有限域的 Galois 群}
$\F_{p^n}/\F_p$ 是可分多项式 $x^{p^n}-x$ 的分裂域. Frobenius 的阶为 $n$, 而域次数也是 $n$, 所以全部自同构恰为其幂, 群为 $C_n$. 固定 $\mathrm{Fr}^d$ 的元素满足 $x^{p^d}=x$; 一般固定域为 $\F_{p^{\gcd(n,d)}}$, 因群中 $\mathrm{Fr}^d$ 生成的子群等于 $\mathrm{Fr}^{\gcd(n,d)}$ 生成的子群. 这恢复上一章的子域分类.
\originalsection{带解答练习}
\begin{exercise}求 $x^4-2$ 的分裂域次数和 Galois 群.\end{exercise}
\begin{solution}
设 $\alpha=\sqrt[4]2$. 艾森斯坦给 $[\Q(\alpha):\Q]=4$, 这是实域. 加入 $i$ 后次数加倍, 得分裂域 $L=\Q(\alpha,i)$ 次数8. 在四根 $\alpha,i\alpha,-\alpha,-i\alpha$ 上, 有自同构 $r$ 送 $\alpha\mapsto i\alpha$, $i\mapsto i$, 因 $L/\Q(i)$ 次数4且 $x^4-2$ 在 $\Q(i)$ 为 $\alpha$ 的最小多项式; 还有复共轭 $s$, 固定 $\alpha$ 并送 $i\mapsto-i$. 它们满足 $r^4=s^2=1,srs=r^{-1}$, 生成8个不同自同构, 所以群为阶8的 $D_4$.
\end{solution}
\begin{remark}
根式可解性与可解群的等价定理需要进一步的根式塔论证. 本册不调用它证明任何结果, 也不把“Galois 群非交换”误当作不可根式求解. 上述三次和四次例子仅计算分裂域和群.
\end{remark}

\originalchapter{多项式恒等式与素性检验}
\originalsection{Fermat 检验的局限}
素数 $p$ 的非零剩余类成阶 $p-1$ 群, 因而 $a^{p-1}\equiv1\pmod p$ 对 $p\nmid a$ 成立. 若某个与 $n$ 互素的 $a$ 不满足该同余, 则 $n$ 合成; 反向不成立. $561=3\cdot11\cdot17$ 对任意整数 $a$ 满足 $a^{561}\equiv a\pmod{561}$: 在每个素因子模下, 若 $a$ 为0结论成立, 否则群阶2,10,16均整除560, 再用中国剩余定理. 这种合数称 Carmichael 数.
\originalsection{多项式素性判别}
\begin{proposition}
设 $n\ge2$ 且 $\gcd(a,n)=1$. 则 $n$ 为素数当且仅当 $(x+a)^n=x^n+a$ 在 $(\Z/n\Z)[x]$ 中成立.
\end{proposition}
\begin{proof}
素数情形由 Frobenius 恒等式及 $a^n=a$. 合数情形取素数 $p\mid n$, 写 $n=p^km$, $p\nmid m$. $x^p$ 的系数为 $\binom np a^{n-p}$. 公式
\[
\binom np=\frac np\binom{n-1}{p-1}
\]
中第二因子模 $p$ 为 $\binom{-1}{p-1}\equiv1$, 因逐项分子为 $-1,-2,\ldots,-(p-1)$ 且分母可逆. 因而该二项系数恰含 $p^{k-1}$, 不被 $p^k$ 整除, $a$ 又不含 $p$, 所以该系数模 $n$ 非零. $p<n$, 右侧无 $x^p$ 项, 恒等式失败.
\end{proof}
直接算此恒等式需要处理 $n$ 个系数, 并未得到关于输入位数的快速算法. AKS 通过模 $x^r-1$ 压缩计算, 其完整参数界与正确性证明不在本册覆盖范围. 本册不复述原讲义未经此处证明的复杂度指数.
\originalsection{平方根见证}
\begin{proposition}
若奇数 $n$ 模下存在 $x\not\equiv\pm1$ 而 $x^2\equiv1$, 则 $n$ 合成, 并可从 $\gcd(x-1,n)$ 得到真因子.
\end{proposition}
\begin{proof}
素数域中 $(x-1)(x+1)=0$ 强迫 $x=\pm1$. 现设 $d=\gcd(x-1,n)$. 若 $d=1$, 则 $x-1$ 可逆, 推得 $x+1=0$ 模 $n$, 矛盾; 若 $d=n$, 则 $x=1$, 也矛盾. 所以 $1<d<n$.
\end{proof}
必须注意反命题错误: 模9的1只有两个平方根1,8, 但9是合数. 一般奇素数幂 $p^k$ 的1也只有两根, 因 $\gcd(x-1,x+1)$ 只可能含因子2, 故整个 $p^k$ 必整除其中一项. 只有含至少两个不同奇素因子的 $n$, 中国剩余定理才给至少四个符号组合根. 原18.703第22讲相关引理的“合数必至少四根”须加此条件.

设 $n\ge3$ 为奇数, 写 $n-1=2^sd$, $d$ 奇, $s\ge1$. 对与 $n$ 互素的底数 $a$, 若 $n$ 素, 序列 $a^d,a^{2d},\ldots,a^{2^sd}=1$ 必或者第一项已为1, 或在首次出现1的前一项为 $-1$. 这是域内平方根仅 $\pm1$ 的结果. 因而若 $a^d\not\equiv1$ 且所有前面平方链项都不为 $-1$, 可判 $n$ 合成. 这是强伪素数检验的确定见证方向; 未证明的错误概率上界不在此宣称.
\originalsection{带解答练习}
\begin{exercise}用平方根见证分解15, 并解释为什么同一办法不能仅靠模9的平方根判合数.\end{exercise}
\begin{solution}
$4^2\equiv1$ 模15但4不是 $\pm1$, $\gcd(4-1,15)=3$ 给真因子. 模9仅1和8满足平方为1, 因而没有非平凡平方根见证; 这不妨碍试除3或其他底数的幂检验判合数.
\end{solution}

\originalchapter{课程作业与中文题解}
本章按18.703 Spring 2013官方11份作业单的原编号收录完整公开题面. 官网未公开官方答案, 本章解答为AI独立撰写并数学交叉审读. 有书目出处但题面完整重印在作业单中的题仍保留出处; 只给外书题号的条目不补造题面. 字母 $C_m$ 表示阶 $m$ 循环群, $D_n$ 阶为 $2n$.

11份作业共126个原题号, 其中45题有完整公开题面, 69个明确小问; 其余81题只给书目定位. 模拟卷另有24题、47小问, 合计69个公开大题号、116小问 (去重前). 作业5第12题与第一次模拟期中第5题相同, 保留双映射; 作业5第9、11题是同一缺面定位, 故81个缺面题号只有80个不同定位. 其余近似主题保留各自原号. 计数把未分问的题视为一个小问, 不把定义中的列举条件另算小问.
\begin{longtable}{@{}p{.12\textwidth}p{.34\textwidth}p{.46\textwidth}@{}}
\toprule 作业&完整题面原题号&仅书目定位, 本册准确登记缺面\\\midrule\endhead
1&无&1--11\\
2&1,3,6&2,4,5,7--14\\
3&1,5,6&2--4,7\\
4&2,5&1,3,4,6--11\\
5&3,10,12--14&1,2,4--9,11\\
6&1,6,8--13&2--5,7\\
7&9&1--8,10--13\\
8&3,5,10&1,2,4,6--9,11--13\\
9&1--12&无\\
10&3,7&1,2,4--6,8\\
11&5--10&1--4\\
\bottomrule
\end{longtable}
逐题书名、章节、节和题号见随包 assessment-map.json; 对应官方原作业单链接可点击. 不拿当今 Judson 的同号题猜2013版题面, 不复制未核许可的 Herstein 外书题面. 作业1全部仅书目定位, 所以没有可诚实重构的题解. 下列章节保留原文件所列截止日期, 不改成新学期日历.
\originalsection{作业2: 2月28日}
\begin{exercise}原第1题. 令 $Z(G)=\{g:hg=gh\text{ 对所有 }h\in G\}$. (i) 证明 $Z(G)=\bigcap_{g\in G}C_G(g)$. (ii) 证明中心为子群.\end{exercise}
\begin{solution}
中心定义恰要求同时属于每个中心化子, 所以等式成立. 中心化子含单位, 且若 $x,y$ 与 $g$ 交换, 则 $xy^{-1}$ 也与 $g$ 交换, 故均为子群, 它们的交也是子群. 亦可直接用同一判别式验证中心.
\end{solution}
\begin{exercise}原第3题, 勘正条件. 非平凡群 $G$ 没有非平凡真子群, 证明它为素数阶循环群. 原题“没有真子群”须作此解释, 否则遗漏平凡群边界.\end{exercise}
\begin{solution}
任取 $a\ne e$, $\langle a\rangle$ 非平凡, 必为整个 $G$. 若无限阶, $\langle a^2\rangle$ 为非平凡真子群; 若有限合数阶 $n$, 取 $1<d<n$, $d\mid n$, 则 $\langle a^d\rangle$ 同样违背条件. 故阶为素数.
\end{solution}
\begin{exercise}原第6题. 找出正方形对称群 $D_4$ 的全部子群.\end{exercise}
\begin{solution}
写 $D_4=\langle r,s\mid r^4=s^2=e,srs=r^{-1}\rangle$. 子群为 $\{e\}$, $\langle r^2\rangle$, 四个 $\langle r^ks\rangle$ ($0\le k<4$), $\langle r\rangle$, $\langle r^2,s\rangle$, $\langle r^2,rs\rangle$, 全群, 共10个. 任意子群与旋转子群交为阶1、2或4的子群. 若还含翻折 $r^ks$, 任意另一个翻折与它的乘积是旋转, 因而子群由该交和这一翻折生成. 对三种交分别枚举正是上述列表, 无遗漏.
\end{solution}
\originalsection{作业3: 3月7日}
\begin{exercise}原第1题, Herstein第3章\S1第1、5题 (题面在作业单完整给出). 求下列乘积并求其阶, 从右向左复合:
\begin{align*}
\text{(i)}\quad&\begin{pmatrix}1&2&3&4&5&6\\6&4&5&2&1&3\end{pmatrix}
\begin{pmatrix}1&2&3&4&5&6\\2&3&4&5&6&1\end{pmatrix};\\
\text{(ii)}\quad&\begin{pmatrix}1&2&3&4&5\\2&1&3&4&5\end{pmatrix}
\begin{pmatrix}1&2&3&4&5\\3&2&1&4&5\end{pmatrix};\\
\text{(iii)}\quad&\alpha^{-1}\begin{pmatrix}1&2&3&4&5\\2&1&3&4&5\end{pmatrix}\alpha,
\quad\alpha=\begin{pmatrix}1&2&3&4&5\\4&1&3&2&5\end{pmatrix}.
\end{align*}
\end{exercise}
\begin{solution}
(i) 逐点先右后左得到 $1\mapsto4\mapsto1$, $2\mapsto5\mapsto3\mapsto2$, 6固定, 为 $(14)(253)$, 阶6. (ii) 为 $(132)$, 阶3. (iii) 换位共轭为 $(\alpha^{-1}(1)\ \alpha^{-1}(2))=(24)$, 阶2.
\end{solution}
\begin{exercise}原第5题. 对 $\sigma=(1472)(365)$, $\tau=(123)(475)$, 求 $\tau\sigma\tau^{-1}$ 及 $\sigma,\tau$ 的阶.\end{exercise}
\begin{solution}
将每个循环位置用 $\tau$ 的像替换, 得 $(2753)(164)$. 两原置换各不交循环长度为4与3, 3与3, 所以阶分别12与3.
\end{solution}
\begin{exercise}原第6题. 找 $\tau\in S_7$ 使 $\tau(125)(3674)\tau^{-1}=(314)(2765)$.\end{exercise}
\begin{solution}
按两循环的列出顺序要求 $1,2,5\mapsto3,1,4$, $3,6,7,4\mapsto2,7,6,5$. 得 $\tau=(132)(45)(67)$, 逐点代入共轭公式可核.
\end{solution}
\originalsection{作业4: 3月14日}
\begin{exercise}原第2题. 证明 $D_4$ 的旋转子群 $H=\{I,R,R^2,R^3\}$ 正规.\end{exercise}
\begin{solution}
$H$ 指数2, 由第2章的指数2判别正规. 也可对生成元检查: 旋转与旋转交换, 翻折将 $R^k$ 共轭为 $R^{-k}$, 均保持 $H$.
\end{solution}
\begin{exercise}原第5题. $G=S_3,H=\{e,(12)\}$. (i) 列所有左陪集. (ii) 列所有右陪集. (iii) 每个左陪集都是某个右陪集吗?\end{exercise}
\begin{solution}
左陪集为 $H,\{(13),(123)\},\{(23),(132)\}$; 右陪集为 $H,\{(13),(132)\},\{(23),(123)\}$. 逐个右乘或左乘 $(12)$ 可核. 含 $(13)$ 的左右两陪集不同, 且每个方向只有一个陪集包含该元素, 所以(iii)否.
\end{solution}
\originalsection{作业5: 3月21日}
\begin{exercise}原第3题, Herstein第2章\S6第3、4题. 设 $N\trianglelefteq G,\bar M\le G/N$, 令 $M=\{a:aN\in\bar M\}$. (i) 证明 $M\le G,N\subseteq M$. (ii) 若 $\bar M$ 正规则 $M$ 正规.\end{exercise}
\begin{solution}
$M$ 是商投影的原像, 非空且若 $a,b\in M$, $(ab^{-1})N=(aN)(bN)^{-1}\in\bar M$, 故成子群. $N$ 的像是单位, 所以包含 $N$. 若 $\bar M$ 正规, $gmg^{-1}$ 的像为 $(gN)(mN)(gN)^{-1}\in\bar M$, 得 $M$ 正规. 原卷把 $M$ 和 $\bar M$ 写成同一字母, 本册区分以免对象混淆.
\end{solution}
\begin{exercise}原第10题, Herstein第2章\S7第6题. 设 $N\trianglelefteq G$. 若 $a$ 阶为有限 $d$, 证明 $aN$ 的阶 $m$ 有限且整除 $d$.\end{exercise}
\begin{solution}
$(aN)^d=N$, 因而有限阶存在. 写 $d=qm+r$, $0\le r<m$, 得 $(aN)^r=N$, 最小性迫使 $r=0$.
\end{solution}
\begin{exercise}原第12题. 两正规子群交平凡, 证明它们逐元素交换.\end{exercise}
\begin{solution}
此题与第一次模拟期中第5题完全相同, 完整交换子论证统一见习题\ref{ex:commute}. 本题原编号保留, 不另计一份不同题面.
\end{solution}
\begin{exercise}原第13题. 证明 $G\cong H'\times K'$ 当且仅当 $G$ 有正规子群 $H\cong H',K\cong K'$, $H\cap K=\{e\}$ 且 $G=H\vee K$ (由两子群生成).\end{exercise}
\begin{solution}
上一题给出 $h,k$ 交换, 所以 $HK$ 是包含两者的子群, 就是生成子群 $H\vee K=G$. 映射 $H\times K\to G$, $(h,k)\mapsto hk$, 是满射同态, 核由 $h=k^{-1}\in H\cap K$ 知平凡. 反向在直积中取两个坐标子群, 它们正规、交平凡且生成全群, 沿同构搬回 $G$ 即得.
\end{solution}
\begin{exercise}原第14题 (挑战). 给有限集合上一种运算, 满足除结合律外的全部群公理.\end{exercise}
\begin{solution}
取五元素 $e,a,b,c,d$, 运算表为
\begin{center}\begin{tabular}{c|ccccc}
$\cdot$&$e$&$a$&$b$&$c$&$d$\\\hline
$e$&$e$&$a$&$b$&$c$&$d$\\
$a$&$a$&$e$&$c$&$d$&$b$\\
$b$&$b$&$d$&$e$&$a$&$c$\\
$c$&$c$&$b$&$d$&$e$&$a$\\
$d$&$d$&$c$&$a$&$b$&$e$
\end{tabular}\end{center}
表定义封闭运算, $e$ 为双侧单位元, 每个元素平方为 $e$ 所以有双侧逆元. 但 $(aa)b=b$, $a(ab)=ac=d$, 故不结合.
\end{solution}
\originalsection{作业6: 4月3日}
\begin{exercise}原第1题. 给二面体群 $D_n$ 的一个生成元关系表示.\end{exercise}
\begin{solution}
对 $n\ge3$, 表示为 $\langle r,s\mid r^n=s^2=e,srs=r^{-1}\rangle$. 关系将任何词整理为 $r^i$ 或 $r^is$, 至多 $2n$ 种, 正 $n$ 边形的具体旋转翻折满足这些关系且有 $2n$ 个不同变换, 所以所表示的群恰为 $D_n$.
\end{solution}
\begin{exercise}\label{ex:3generate}原第6题, Herstein第3章\S3第6题. 证明3循环生成 $A_n$.\end{exercise}
\begin{solution}
偶置换可按相邻两换位配对. 相同两换位抵消; 共一个字母的两换位之积是3循环; 不交情形有 $(ab)(cd)=(acb)(acd)$, 从右向左作用可逐点核实. 所以每对均可化为3循环的积, 而每个3循环本来是偶置换.
\end{solution}
\begin{exercise}原第8题. 证明 Sylow $p$ 子群为特征子群当且仅当它唯一.\end{exercise}
\begin{solution}
特征子群指被每个自同构保持, 特别被内自同构保持而正规, Sylow 共轭性给唯一. 反之自同构保持子群阶, 必置换所有 Sylow $p$ 子群; 唯一时就必须保持它.
\end{solution}
\begin{exercise}原第9题. 交换群 $G$ 阶为 $p^lq^m$, $p\ne q$, $l,m>0$. 证明 $G\cong P\times Q$, $|P|=p^l,|Q|=q^m$.\end{exercise}
\begin{solution}
取两 Sylow 子群, 交换性使它们正规, 阶互素使交平凡. 乘积子群阶 $|P||Q|=|G|$, 故生成全群. 作业5第13题的同构 $(a,b)\mapsto ab$ 给直积.
\end{solution}
\begin{exercise}原第10题. 证明阶168的单群同构于 $S_8$ 的子群.\end{exercise}
\begin{solution}
$n_7\mid24$, $n_7\equiv1\pmod7$, 只能1或8. 单性排除1. 共轭作用在八个 Sylow 7子群上是非平凡同态 $G\to S_8$, 核正规且不是全群, 故平凡, 得嵌入.
\end{solution}
\begin{exercise}\label{ex:an}原第12题 (挑战), Herstein第3章\S3第9、10题. (i) $n\ge5$, 若 $\{e\}\ne N\trianglelefteq A_n$, 证明 $N$ 含3循环. (ii) 在 $n\ge3$ 范围, 证明 $A_n$ 单当且仅当 $n\ne4$. 原题没有注明(ii)的低阶范围, $A_1,A_2$ 平凡不算单群.\end{exercise}
\begin{solution}
取 $\sigma\in N\setminus\{e\}$ 使移动点数 $m$ 最小. 若 $m\ge7$, 取 $a$ 使 $\sigma(a)=b\ne a$, 选不同的 $c,d$ 避开 $a,b$, 令 $\tau=(acd)$. $\sigma\tau\sigma^{-1}$ 的支撑含 $b$, $\tau$ 的支撑不含 $b$, 故不相等. 交换子 $\sigma\tau\sigma^{-1}\tau^{-1}\in N$ 非平凡, 最多移动两个三点支撑的并, 即6点, 矛盾. 所以 $3\le m\le6$.

$m=3$ 即3循环. $m=4$ 必为 $(ab)(cd)$; 选支撑外一点 $e_0$, 用 $\tau=(cde_0)$, 共轭为 $\tau^{-1}$, 交换子为 $\tau^{-2}=\tau$, 得3循环. $m=5$ 必为 $(abcde_0)$, 用 $\tau=(abc)$, 交换子 $(bcd)(acb)=(adb)$ 是3循环. $m=6$ 的偶置换类型仅 $(abcd)(e_0f)$ 或 $(abc)(de_0f)$. 前者同样取 $(abc)$ 得 $(adb)$; 后者取 $\tau=(abd)$, 共轭为 $(bce_0)$, 二者支撑不同, 交换子非平凡而至多移动5点, 与最小性矛盾. 这穷尽了无固定点的偶循环类型, 证明(i).

任意两个3循环在 $A_n$ 中共轭: 先在 $S_n$ 中逐位置匹配三点; 若匹配置换为奇, 再在目标循环补集两个点交换, 因 $n\ge5$ 存在这两点, 改为偶且保持匹配. 正规性使 $N$ 含全部3循环, 由习题\ref{ex:3generate}生成全 $A_n$, 得单性. $A_3=C_3$ 单, $A_4$ 的双换位子群 $V$ 非平凡真且正规, 故不单. 这给准确低阶边界.
\end{solution}
\begin{exercise}原第11题与第13题 (后者挑战). 分别证明阶在1至59、61至167的单群必为素数阶.\end{exercise}
\begin{solution}
先给统一约束. 交换单群必为素数阶循环群, 因每个子群正规且任意非单位元生成全群. 非交换单群 $G$ 对任何真子群 $H$ 的左陪集作用非平凡, 核正规故忠实. 符号同态必平凡, 否则 $G$ 嵌入阶2群; 因而 $G\hookrightarrow A_m$, $m=[G:H]$, 必须 $|G|\mid m!/2$. 同样, 若 Sylow 数 $n_p>1$, 共轭作用给 $G\hookrightarrow A_{n_p}$. 合数素数幂阶由中心非平凡排除: 单性使中心为全群, 会迫使交换, 只能素数阶.

下表对每个列出合数 $n$, 取 Sylow $p$ 子群 $H$, $m=n/p^{v_p(n)}$, 已直接核 $n\nmid m!/2$, 故排除. $v_p(n)$ 是整数分解中 $p$ 的指数. 表格是有限算术, 作用约束才是数学理由.
\begin{longtable}{@{}rrr@{\qquad}rrr@{\qquad}rrr@{}}
\toprule $n$&$p$&$m$&$n$&$p$&$m$&$n$&$p$&$m$\\\midrule\endhead
6&2&3&10&5&2&12&2&3\\
14&7&2&15&5&3&18&3&2\\
20&5&4&21&7&3&22&11&2\\
24&2&3&26&13&2&28&7&4\\
33&11&3&34&17&2&35&7&5\\
36&3&4&38&19&2&39&13&3\\
40&2&5&42&7&6&44&11&4\\
45&3&5&46&23&2&48&2&3\\
50&5&2&51&17&3&52&13&4\\
54&3&2&55&11&5&57&19&3\\
58&29&2&&&&&&\\
\bottomrule
\end{longtable}
1至59中余下合数仅30和56. 阶30单性迫使 $n_5=6,n_3=10$, 产生 $6\cdot4+10\cdot2=44$ 个不同非单位元, 超过29. 阶56单性迫使 $n_7=8$, 产生48个阶7元素; 剩余恰8元素, 任何阶8 Sylow 2子群都须等于此剩余集合, 因而唯一正规, 矛盾. 第11题得证.

61至167的陪集排除表如下:
\begin{longtable}{@{}rrr@{\qquad}rrr@{\qquad}rrr@{}}
\toprule $n$&$p$&$m$&$n$&$p$&$m$&$n$&$p$&$m$\\\midrule\endhead
62&31&2&65&13&5&66&11&6\\
68&17&4&69&23&3&74&37&2\\
75&5&3&76&19&4&77&11&7\\
78&13&6&80&2&5&82&41&2\\
85&17&5&86&43&2&87&29&3\\
88&11&8&91&13&7&92&23&4\\
93&31&3&94&47&2&95&19&5\\
96&2&3&98&7&2&99&11&9\\
100&5&4&102&17&6&104&13&8\\
106&53&2&108&3&4&110&11&10\\
111&37&3&112&2&7&114&19&6\\
115&23&5&116&29&4&117&13&9\\
118&59&2&119&17&7&122&61&2\\
123&41&3&124&31&4&129&43&3\\
130&13&10&133&19&7&134&67&2\\
135&3&5&136&17&8&138&23&6\\
141&47&3&142&71&2&143&13&11\\
145&29&5&146&73&2&147&7&3\\
148&37&4&150&5&6&152&19&8\\
153&17&9&155&31&5&156&13&12\\
158&79&2&159&53&3&160&2&5\\
161&23&7&162&3&2&164&41&4\\
166&83&2&&&&&&\\
\bottomrule
\end{longtable}
阶63,70,84,126,140,154,165分别取素数7,5,7,7,5,11,11, Sylow 整除和同余条件只允许 $n_p=1$, 所以不单. 阶72时 $n_3=1$ 或4; 后者迫使嵌入 $A_4$, 阶不容. 余下仅90,105,120,132,144, 逐一处理.

阶90: 单性给 $n_3=10,n_5=6$. 阶9 Sylow 子群交换. 若不同两者 $Q,Q'$ 交非平凡, 取交中的阶3元 $x$, 则两者都在 $C_G(x)$ 中. 此中心化子阶为9,18,45或90. 前三种阶的 Sylow 3子群都唯一 (个数整除1,2,5并模3余1); 最后一种使 $x$ 中央, 与非交换单性冲突. 所以两者交平凡. 产生 $10\cdot8+6\cdot4=104$ 个不同非单位元, 超过89.

阶105: 单性给 $n_7=15,n_5=21$, 两类阶素数元素共 $15\cdot6+21\cdot4=174$, 超过104. 阶120的完整排除见第二次模拟期中习题\ref{ex:120}; 其中所需 $A_6$ 单性已在本章前题证明. 阶132: $n_{11}=12$, $n_3=1,4$ 或22, 前两值被正规性或嵌入 $A_4$ 排除. 后值产生 $12\cdot10+22\cdot2=164$ 个不同非单位元, 超过131.

阶144: $n_3=1,4$ 或16, 单性和 $A_4$ 阶限制迫使16. 若不同阶9子群交非平凡, 对交中的阶3元 $x$, 中心化子阶为9,18,36,72或144. 前两阶只有一个 Sylow 3子群; 36,72的中心化子指数4,2, 陪集作用的忠实性不容阶144的单群; 144使 $x$ 中央. 所以16个 Sylow 3子群只交于单位, 给128个非单位奇数阶元素. 剩余恰16元素, 每个阶16 Sylow 2子群都等于剩余集合, 成为唯一正规子群, 矛盾. 两表、素数幂和这些余项覆盖区间全部合数, 第13题得证.
\end{solution}
\originalsection{作业7: 4月11日}
\begin{exercise}原第9题, Herstein第4章\S1第26题. 复四元数 $H(\C)$ 的元素为 $\alpha_0+\alpha_1i+\alpha_2j+\alpha_3k$, 系数在 $\C$, 四元数基元满足 $i^2=j^2=k^2=ijk=-1$. 证明它不是除环.\end{exercise}
\begin{solution}
用 $\iota\in\C$ 表示标量虚数, 避免与基元 $i$ 混淆. 两个非零元素 $1+\iota i$ 与 $1-\iota i$ 的乘积是 $1-(\iota i)^2=1-1=0$, 因标量与基元交换. 有非零零因子, 不可能每个非零元可逆: 若第一因子可逆, 左乘其逆会迫使第二因子为0.
\end{solution}
\originalsection{作业8: 4月18日}
\begin{exercise}原第3题. 对理想 $I,J$ 证明 (i) $I\cap J$ 为理想; (ii) $I+J=\{i+j\}$ 为理想; (iii) 有限和 $\sum i_kj_k$ 组成的 $IJ$ 为理想.\end{exercise}
\begin{solution}
交中元素的差以及左右乘法仍在两理想中. 和中 $(i+j)-(i'+j')=(i-i')+(j-j')$ 留在和中, 左右乘 $r$ 分别用 $ri,rj$ 或 $ir,jr$ 留在原理想. 积中和与负元逐项构造, 左乘将每项变成 $(ri_k)j_k$, 右乘变成 $i_k(j_kr)$, 仍为允许有限和. 这些证明对非交换环的双边理想也成立.
\end{solution}
\begin{exercise}原第5题. 固定素数 $p$, 令 $R$ 为既约分母不被 $p$ 整除的有理数. (i) 证明 $R\le\Q$ 为子环. (ii) 分子被 $p$ 整除的元素组成理想 $I$, 证明 $R/I\cong\F_p$.\end{exercise}
\begin{solution}
加乘运算后的分母可取原两分母的积, 不被 $p$ 整除, 约分也不引入 $p$, 且0,1及负元均在集合中. 映射 $a/b\mapsto\bar a\bar b^{-1}$ 良定义: 相等分数有 $ad=bc$, 在 $\F_p$ 中消去可逆分母. 映射保加乘及1, 整数像使它满射, 核恰是分子被 $p$ 整除的集合, 因而为理想, 第一同构定理给结论.
\end{solution}
\begin{exercise}原第10题, Herstein第4章\S3第22、23题. 互素正整数 $m,n$, 令 $I_m=m\Z,I_n=n\Z$. (i) 求交. (ii) 建立 $\Z/(mn)\to\Z/(m)\times\Z/(n)$ 的单射环同态. (iii) 证明它为同构.\end{exercise}
\begin{solution}
(i) 同时整除条件和互素性给交 $mn\Z$. (ii) 映射 $\bar a\mapsto(a\bmod m,a\bmod n)$ 保加乘, 良定义, 核由(i)为0. (iii) 取 $um+vn=1$, 对给定余数 $A,B$, $Avn+Bum$ 同时满足模 $m$ 为 $A$, 模 $n$ 为 $B$, 得满射. 若原模数为负, 理想不变, 改取绝对值; 模1允许零商环例外.
\end{solution}
原作业末尾还要求证明环直积有一般余积泛性质, 这个要求本身错误. 第5章练习以两恒等映射给出反例; 本册登记此附注并纠正, 不为假命题伪造证明.
\originalsection{作业9: 4月25日}
\begin{exercise}原第1题, Herstein第4章\S4第2题. 对高斯整数环, 分量均被3整除的集合 $M$ 为理想, 证明商为9元域.\end{exercise}
\begin{solution}
$M=(3)$, 所以是理想. 商同构于 $\F_3[x]/(x^2+1)$, 对应 $x\mapsto i$. 每个剩余类恰由两个模3分量确定, 共9个. $x^2+1$ 在 $\F_3$ 无根 (平方仅0,1), 故不可约, 商为域.
\end{solution}
\begin{exercise}原第2题, Herstein第4章\S4第3题. (i) 证明 $R=\Z[\sqrt2]$ 为复数子环. (ii) 分量均被5整除的 $M$ 为理想, 且商为25元域.\end{exercise}
\begin{solution}
乘积 $(a+b\sqrt2)(c+d\sqrt2)=(ac+2bd)+(ad+bc)\sqrt2$ 留在集合中, 加负元及1也成立. $M=(5)$, 商同构于 $\F_5[x]/(x^2-2)$, 共25个元素. 模5平方为0,1,4, 不含2, 二次式不可约, 故商域成立. 也可由共轭乘积 $a^2-2b^2$ 给非零类求逆.
\end{solution}
\begin{exercise}原第3题. 构造49元域.\end{exercise}
\begin{solution}
取 $\F_7[x]/(x^2+1)$. 模7平方为0,1,2,4, 不含 $-1=6$, 所以二次式不可约; 商域基为 $1,\bar x$, 有 $7^2=49$ 个元素.
\end{solution}
\begin{exercise}原第4题. 真理想 $I$ 外的每个元素均为单位, 证明 $I$ 为唯一极大理想.\end{exercise}
\begin{solution}
任一真理想不能含单位, 否则吸收其逆元得到1. 因此任何真理想均包含于 $I$, 其中 $I$ 本身为真理想, 所以既极大又唯一.
\end{solution}
\begin{exercise}原第5题. 环同态 $\varphi:R\to S$ 保1. (i) 素理想 $J\subset S$ 的原像为素理想. (ii) 给 $J$ 极大而原像非极大的例子.\end{exercise}
\begin{solution}
核验差和标量吸收即得原像为理想; 保1使1不在原像, 所以为真理想. 若 $ab$ 在原像, $\varphi(a)\varphi(b)\in J$, 素性给一因子在 $J$, 从而一因子在原像. 对(ii), 用 $\Z\hookrightarrow\Q$ 和 $J=(0)$, 域中的零理想极大, 原像为 $\Z$ 的零理想, 却严格位于 $(2)$ 下, 非极大.
\end{solution}
\begin{exercise}原第6题. 整环中证明 (i) $a\mid b\iff(b)\subseteq(a)$; (ii) $a,b$ 相伴当且仅当主理想相同; (iii) $a$ 为单位当且仅当 $(a)=R$.\end{exercise}
\begin{solution}
(i) $b=ac$ 使所有 $rb$ 是 $a$ 的倍数, 反向包含使 $b$ 本身为倍数. (ii) 相伴直接给相同理想; 反向写 $a=ub,b=va$, 若非零消去得 $uv=1$, 因交换环 $u$ 单位. 若一方为0, 理想相同迫使另一方也0, 二者相伴. (iii) $(a)=R$ 等价于 $1=ra$, 正是可逆性.
\end{solution}
\begin{exercise}原第7题. 证明整环的素元不可约.\end{exercise}
\begin{solution}
若素元 $p=ab$, 素性使 $p\mid a$ 或 $p\mid b$. 如 $a=pc$, 消去非零 $p$ 得 $1=cb$, 所以 $b$ 单位. 另一情况对称.
\end{solution}
\begin{exercise}原第8题. 在整环中, 同一对元素的两个 gcd 必相伴.\end{exercise}
\begin{solution}
每个 gcd 既是公因子又被所有公因子整除, 故二者互相整除. 第6题(ii)给相伴, 包括零情形.
\end{solution}
\begin{exercise}原第9题. 对 $R=\Z[\sqrt{-5}]$: (i) 证明 $2,3,1\pm\sqrt{-5}$ 不可约. (ii) 证明每个非零非单位有不可约分解. (iii) 证明非UFD.\end{exercise}
\begin{solution}
范数 $a^2+5b^2$ 乘法性成立, 单位恰 $\pm1$, 且范数2,3均无解. 四元素范数为4,9,6,6, 非平凡因子分解会要求范数2或3, 所以不可约. 任意可约非单位的两个非单位因子的范数均严格小于原范数, 按正整数范数归纳给有限分解. 最后 $6=2\cdot3=(1+\sqrt{-5})(1-\sqrt{-5})$, 不可约因子范数不同故无法相伴, 唯一性失败.
\end{solution}
\begin{exercise}原第10题. UFD中 (i) 证明 lcm 存在, 每个公倍数都被它整除, 且仅差相伴; (ii) 证明 gcd 与 lcm 的乘积与 $ab$ 相伴.\end{exercise}
\begin{solution}
对非零 $a,b$ 按全部不可约相伴类写指数 $r_i,s_i$, lcm 取最大指数, gcd 取最小指数. 整除等价于逐指数比较, 给存在、定义条件及唯一性. 因 $\min(r_i,s_i)+\max(r_i,s_i)=r_i+s_i$, 乘积与 $ab$ 相伴. 若一个为0, lcm为0、gcd取另一元; 两者都0时均取0, 条件仍成立.
\end{solution}
\begin{exercise}原第11题 (挑战). 在有单位元的交换半群中定义唯一分解.\end{exercise}
\begin{solution}
原题的“半群”实际含单位元, 本册称交换幺半群. 单位 $u$ 指存在 $v$ 使 $uv=1$. 相伴指相差一个单位因子. 非单位 $a$ 为原子指 $a=bc$ 迫使某因子为单位. 唯一分解指每个非单位为有限原子乘积, 且任意两个原子分解的原子数相同, 经重排后逐项相伴. 单位以空积和一个单位系数处理. 一般不消去半群可能有异常, 不能无条件套用整环证明; 下一题在整数向量子幺半群中可消去.
\end{solution}
\begin{exercise}原第12题 (挑战). 给 $v_1,\ldots,v_n\in\Z^2$, 令 $S$ 是其非负整数线性组合, 运算为加法. 确定哪些 $S$ 有唯一分解.\end{exercise}
\begin{solution}
必须先处理单位: $U=S\cap(-S)$ 是单位群, $s,t$ 相伴当且仅当 $s-t\in U$. 在约化幺半群 $\bar S=S/U$ 中, 取全部不同原子类 $\bar a_1,\ldots,\bar a_r$. 精确判据为这些类在群 $\operatorname{gp}(S)/U$ 中无非零整数关系; 等价于映射 $\N^r\to\bar S$, $(n_i)\mapsto\sum n_i\bar a_i$, 为双射. $\operatorname{gp}(S)$ 表示生成的整数加法群.

说明原子确实有限且生成. 令 $C$ 为原生成向量的非负实锥, $L=C\cap(-C)$ 为其直线部分. 有 $L=\operatorname{span}_{\R}U$. 一方面单位方向显然在 $L$; 另一方面, 位于 $L$ 的生成向量的负方向可由原生成向量作非负有理组合 (整数系数线性系统存在实非负解时可取有理解). 清分母使 $-kv\in S$, 加上 $(k-1)v$ 得 $-v\in S$, 该生成向量为单位. 直线部分由落在其中的生成向量张成: 若 $x,-x\in C$, 将两表达相加为0, 参与该零和的生成方向均有负向在锥中. 因此两边实张成相同.

在 $\R^2/L$ 中, 非单位生成方向组成有限尖锥, 可选有理线性函数 $\ell$, 对它们严格正且在 $L$ 为0. 一维时它们同向; 二维时尖锥的两个极端方向夹角小于半周, 取其内侧法向之间的有理方向即可使内积全正. 清分母使 $\ell(S)$ 为非负整数, 且非单位取正值. 若非单位分成两个非单位, 两者次数 $\ell$ 均下降, 按次数归纳得到有限原子分解. 原子在原生成向量表达中只能有一个非单位生成元, 否则可拆, 其余都是单位. 所以原子类来自有限原列表, 删去单位、重复伴随类及可分解者后得到全部原子类.

这些类生成 $\bar S$, 满射性已有. 两个非负指数组合相同恰是两个不同原子分解, 所以单射性恰为唯一性. 非零整数关系拆成正负两边便产生不同分解; 它不能只含同号项, 因非单位之和不可能是单位: 若 $a+b+c=0$, 则 $b+c$ 是 $a$ 的逆, 使 $a$ 已是单位.

向量和 $U$ 的基为整数, 在商实空间的线性关系可取有理系数, 清分母并再乘一个整数使落入 $U$ 的格, 故整数无关系等价于 $a_i$ 在 $\R^2/\operatorname{span}_{\R}U$ 中线性无关. 因而分类为: $\operatorname{rank}U=0$ 时, $S=\{0\}$ 或由一个非零向量或两个不共线向量自由生成; $\operatorname{rank}U=1$ 时, $S=U$ 或 $S=U+\N a$ 且 $a\notin\operatorname{span}_{\R}U$; $\operatorname{rank}U=2$ 时 $S=U$, 全是单位. 例如 $(2,0),(0,3)$ 生成的幺半群唯一分解, 无须是整个 $\Z^2$ 的格基; $(2,0),(3,0)$ 则有 $3(2,0)=2(3,0)$, 不唯一.
\end{solution}
\originalsection{作业10: 5月2日}
\begin{exercise}原第3题. 求 $\Z[i]$ 中 $11+7i$ 与 $8-i$ 的 gcd.\end{exercise}
\begin{solution}
作欧几里得除法 $11+7i=(1+i)(8-i)+2$, 再 $8-i=4\cdot2-i$. 余项 $-i$ 为单位, 故 gcd 与1相伴. 可复核的贝祖式为 $1=-4i(11+7i)+(-4+5i)(8-i)$.
\end{solution}
\begin{exercise}原第7题 (挑战). 给一个环和元素 $a$, 它存在不可约分解, 但从 $a$ 开始的因子分解算法有不停止的选择路径.\end{exercise}
\begin{solution}
取独立变量 $u,v,w$ 和 $A=\Q[u,v,w,w^{-1}]$, 定义子环
\[
R=\Q[u,v,w,uv/w,uv/w^2,uv/w^3,\ldots].
\]
它为整环. 其中单项式 $u^av^bw^c$ 的指数 $a,b\ge0$, 且 $c<0$ 时必须 $a,b\ge1$; 反向每个这样的单项式都由所列生成元得到, 这是检查因子属于 $R$ 的依据.

$A$ 是将 UFD $\Q[u,v,w]$ 中的 $w$ 变为单位得到的环, 仍唯一分解: 清掉 Laurent 分母后在多项式环中分解, 只把 $w$ 的幂并入单位. 其单位恰为 $\lambda w^j$, $\lambda\in\Q^\times,j\in\Z$, 因清分母后的可逆性只容 $w$ 因子. 所以 $u=fg$ 在 $A$ 中, 除交换次序外必为 $f=\lambda uw^j,g=\lambda^{-1}w^{-j}$. 两者若属于 $R$, 后者迫使 $j\le0$, 前者因无 $v$ 而迫使 $j\ge0$, 故 $j=0$, 某因子单位. 因而 $u$ 是 $R$ 原子, $v$ 同理, $a=uv$ 有原子分解.

另一方面令 $a_k=uv/w^k$, 则 $a_k=w a_{k+1}$ 对所有 $k\ge0$ 成立. 所有因子属于 $R$ 且非单位: $a_k$ 在 $A$ 已非单位, $w$ 的逆 $w^{-1}$ 不属于 $R$. 永远选 $a_{k+1}$ 继续拆便无穷, 每步都是非单位分解. 这不否认 $a_0$ 存在另一条有限分解, 也不宣称 $R$ 每个元素都有原子分解.
\end{solution}
\originalsection{作业11: 5月9日}
官网此文件是第11份作业, 原 PDF 标题误写“作业9”; 以下按官方目录顺序标11, 保留日期对应.
\begin{exercise}原第5题. $F[x]$ 自同构 $\varphi$ 固定每个 $a\in F$. (i) 证明 $f$ 不可约当且仅当 $\varphi(f)$ 不可约. (ii) 证明次数保持.\end{exercise}
\begin{solution}
(i) 自同构双向保持乘积和单位, 把非平凡分解双向搬运, 故保持不可约. (ii) 令 $h=\varphi(x),k=\varphi^{-1}(x)$, 它们都非恒定, 否则逆映射固定常数会使 $x$ 为常数. $k(h(x))=x$, 域上非恒定多项式复合次数为次数之积, 得 $\deg h=\deg k=1$. 因 $\varphi$ 固定系数, $\varphi(f)=f(h)$, 对非零 $f$ 次数保持; 零多项式双方同为0, 其次数按不定义或 $-\infty$ 的一致约定处理.
\end{solution}
\begin{exercise}原第6题. $b\in F^\times,c\in F$, 证明 $f(x)\mapsto f(bx+c)$ 为固定 $F$ 的自同构.\end{exercise}
\begin{solution}
代入保持和、积及常数, 逆代入为 $x\mapsto b^{-1}(x-c)$, 两次复合均为 $x$, 所以是自同构.
\end{solution}
\begin{exercise}原第7题. 证明所有固定 $F$ 的 $F[x]$ 自同构都是上述形式.\end{exercise}
\begin{solution}
由第5题(ii), $\varphi(x)=bx+c$ 且 $b\ne0$. 固定每个系数给 $\varphi(\sum a_ix^i)=\sum a_i(bx+c)^i$, 完成分类.
\end{solution}
\begin{exercise}原第8题. (i) 给 $\Q[x]$ 非恒等且平方为恒等的自同构. (ii) 对任意 $n>0$, 给 $\C[x]$ 阶为 $n$ 的自同构.\end{exercise}
\begin{solution}
(i) 取 $x\mapsto-x$. (ii) 取 $x\mapsto\zeta_n x$, $\zeta_n=\exp(2\pi i/n)$ 是本原 $n$ 次单位根, 固定复系数. 第 $k$ 次幂送 $x$ 至 $\zeta_n^kx$, 恰当 $n\mid k$ 才恒等, 所以阶为 $n$, 包括 $n=1$ 的恒等映射.
\end{solution}
\begin{exercise}原第9题. 对特征 $p>0$ 的域: (i) 证明 $(a+b)^p=a^p+b^p$. (ii) 证明 $q=p^n$ 次幂映射为环同态. (iii) 证明单射. (iv) 有限域时为自同构.\end{exercise}
\begin{solution}
特征 $p$ 必素, 二项式中间系数被 $p$ 整除, 给(i). 迭代 $n$ 次得 $q$ 次幂加法性, 乘法与1保持直接成立. 若 $a^q=b^q$, 则 $(a-b)^q=0$, 域无非零幂零元, 得 $a=b$. 有限集单射为满射, 故(iv).
\end{solution}
\begin{exercise}原第10题 (挑战). 给特征 $p$ 的域使 $a\mapsto a^p$ 不满射.\end{exercise}
\begin{solution}
取 $F=\F_p(t)$. 若 $t=(f/g)^p$, 清分母得 $f^p=tg^p$, 不可约因子 $t$ 的指数左边是 $p$ 倍数, 右边是 $1$ 加 $p$ 倍数, 不可能. 因而 $t$ 没有原像.
\end{solution}

\originalchapter{课程模拟考试与中文题解}
本部分题面来自 MIT 18.703, Spring 2013 的3份公开模拟试卷, 原编号、分值及小问保持对应. 官网没有随卷提供答案, 以下解答均为本整理本独立撰写并交叉审读, 不是官方答案. 原卷中的空白答题页合并, 不影印英文整页. 正式两次期中及期末试卷未在此课程考试页面公开.
\originalsection{第一次模拟期中}
原卷限时80分钟, 闭卷、不可使用计算器, 共7题; 题目95分, 表达5分, 合计100分.
\begin{exercise}\label{ex:m1q1}原卷第1题 (15分). (i) 定义正规子群. (ii) 定义群同态. (iii) 定义交错群 $A_n$.\end{exercise}
\begin{solution}
(i) $N\le G$ 且任意 $g$ 有 $gNg^{-1}=N$. (ii) $\varphi:G\to H$ 满足 $\varphi(ab)=\varphi(a)\varphi(b)$. (iii) $A_n$ 是 $S_n$ 中偶置换组成的子群, 即符号同态的核; 偶性由换位分解的偶数长度定义, 其不依赖分解已在第3章证明.
\end{solution}
\begin{exercise}原卷第2题 (15分). (i) 给 $A_4$ 的一个非平凡真正规子群 $V$, 说明其同构类型. (ii) 写其全部左陪集. (iii) 求 $A_4/V$ 的同构类型.\end{exercise}
\begin{solution}
取 $V=\{e,(12)(34),(13)(24),(14)(23)\}\cong C_2\times C_2$, 共轭保留双换位, 故正规. 全部左陪集为
\begin{align*}
V,\quad &(123)V=\{(123),(134),(243),(142)\},\\
&(132)V=\{(132),(234),(124),(143)\}.
\end{align*}
各集合含4个元素且不交, 共覆盖12个元素. 商阶3, 所以为 $C_3$.
\end{solution}
\begin{exercise}原卷第3题 (20分). 对 $H\le G$, 证明下列条件等价: (1) $H$ 正规; (2) 每个 $g$ 有 $gHg^{-1}=H$; (3) 每个 $a$ 有 $aH=Ha$; (4) 全部左陪集组成的集合与全部右陪集组成的集合相同.\end{exercise}
\begin{solution}
(1)与(2)为定义, (2)右乘 $g$ 等价于(3), (3)推出(4). 若(4)成立, 对任一 $a$, 存在 $b$ 使 $aH=Hb$. 因 $a\in Hb$, 有 $a=hb$, $h\in H$, 故 $Ha=Hhb=Hb=aH$, 得(3). 最后一步不能仅从两边“陪集数相同”推出, 必须用到同一陪集中包含 $a$.
\end{solution}
\begin{exercise}原卷第4题 (10分). 设 $N\trianglelefteq G$. 证明 $G/N$ 交换当且仅当 $N$ 包含所有交换子 $aba^{-1}b^{-1}$.\end{exercise}
\begin{solution}
商中 $aN,bN$ 交换当且仅当其交换子陪集为 $N$, 即 $aba^{-1}b^{-1}\in N$. 对所有 $a,b$ 量化, 即得两方向.
\end{solution}
\begin{exercise}\label{ex:commute}原卷第5题 (10分), 与作业5第12题同题. 若 $H,K\trianglelefteq G$ 且交为单位子群, 证明其元素两两交换.\end{exercise}
\begin{solution}
对 $h\in H,k\in K$, $hkh^{-1}k^{-1}$ 由 $K$ 正规属于 $K$, 由 $H$ 正规又属于 $H$. 交平凡使该交换子为 $e$, 即 $hk=kh$.
\end{solution}
\begin{exercise}\label{ex:sylowstate}原卷第6题 (15分). (i) 陈述 Sylow 定理. (ii) 证明不同素数 $p,q$ 的乘积阶群不为单群.\end{exercise}
\begin{solution}
(i) 若 $|G|=p^am$, $p\nmid m$, 则有阶 $p^a$ 子群; 所有这样的子群共轭, 任意 $p$ 子群包含在其中一个内; 个数整除 $m$ 且模 $p$ 余1. 完整证明见定理\ref{thm:sylow}. (ii) 不妨 $p<q$, $n_q\mid p$, $n_q\equiv1\pmod q$, 所以 $n_q=1$. 唯一 Sylow $q$ 子群非平凡、真且正规, 排除单性.
\end{solution}
\begin{exercise}原卷第7题 (10分). 建立包含 $N\trianglelefteq G$ 的子群与 $G/N$ 子群的自然双射, 并证明正规性保持.\end{exercise}
\begin{solution}
对投影 $\pi$, 用 $H\mapsto\pi(H)$ 与 $\bar H\mapsto\pi^{-1}(\bar H)$. 若 $N\subseteq H$, 则 $\pi^{-1}(\pi(H))=H$, 因像相同的两元素相差核元素; 满射性给反向复合为恒等. 子群封闭性由像、原像保持乘法和逆元得出. $\pi(gHg^{-1})=\pi(g)\pi(H)\pi(g)^{-1}$ 且原像唯一, 故正规性双向对应.
\end{solution}
\originalsection{第二次模拟期中}
原卷限时80分钟, 闭卷、不可使用计算器, 共6题; 题目95分, 表达5分, 合计100分.
\begin{exercise}原卷第1题 (15分). (i) 定义整环的不可约元. (ii) 定义素元. (iii) 定义主理想整环.\end{exercise}
\begin{solution}
不可约元是非零非单位, 任一因子分解都含单位因子. 素元是非零非单位, 若整除乘积则整除至少一个因子. PID 是每个理想均由一个元素生成的整环. “整环”及“非单位”条件均不能省略.
\end{solution}
\begin{exercise}\label{ex:120}原卷第2题 (15分). (i) 陈述 Sylow 定理. (ii) 证明不存在阶120的单群.\end{exercise}
\begin{solution}
(i) 见习题\ref{ex:sylowstate}的完整陈述. (ii) 设 $G$ 单, 则 $n_5\mid24$, $n_5\equiv1\pmod5$, 只能为1或6, 单性排除1. 共轭作用在六个 Sylow 5子群上给非平凡同态 $G\to S_6$, 核正规故单射. 符号同态在 $G$ 上也必平凡 (否则单射到阶2群), 所以像在 $A_6$. 于是 $G$ 在 $A_6$ 中指数3. $A_6$ 在其三个左陪集上作用, 得非平凡 $A_6\to S_3$. 作业6第12题的证明给 $A_6$ 单, 迫使此映射单射, 但 $360>6$, 矛盾.
\end{solution}
\begin{exercise}原卷第3题 (15分). 对整环 $R$ 的真理想 $I$, 证明 $R/I$ 是域当且仅当 $I$ 极大.\end{exercise}
\begin{solution}
事实上只需 $R$ 交换有1. 若 $I$ 极大, 对 $a\notin I$, $I+(a)=R$, 所以 $i+ra=1$ 给商中逆元. 若商为域, 任意 $J\supsetneq I$ 中的非零像可逆, 故 $J/I$ 含1, 于是 $J=R$. 真理想条件排除零商环.
\end{solution}
\begin{exercise}原卷第4题 (15分). (i) 证明理想升链的并是理想. (ii) 证明 PID 的理想升链最终稳定.\end{exercise}
\begin{solution}
若 $x,y$ 属并, 它们同属较后一个理想, 故差也在并; 任意环元素乘 $x$ 仍在含 $x$ 的理想. 若并 $I=(d)$, 元素 $d$ 已在某个 $I_N$, 故 $I\subseteq I_N$, 而反向包含本来成立. 所有后续理想都等于 $I_N$.
\end{solution}
\begin{exercise}原卷第5题 (20分). (i) 证明两理想的交是理想. (ii) 并集必为理想吗?\end{exercise}
\begin{solution}
交中的差以及与任意环元素的左右乘积同时在两个理想中, 故在交中. 并不必是理想, 例如 $2\Z\cup3\Z$ 含2和3却不含5. 一般并是加法子群当且仅当一个理想包含另一个: 若各取一个不属于对方的元素, 它们的和不属于任何一方.
\end{solution}
\begin{exercise}原卷第6题 (15分). 在 PID 中证明非零 $a,b$ 有 gcd $d$ 且 $d=ra+sb$.\end{exercise}
\begin{solution}
写理想 $(a,b)=(d)$. 因 $a,b\in(d)$, $d$ 整除二者. 生成元 $d$ 属 $(a,b)$ 故具有所述表达, 任意公因子由表达式整除 $d$. 这就是 gcd 定义, 唯一性仅差单位.
\end{solution}
\originalsection{模拟期末}
原卷限时3小时, 闭卷、不可使用计算器, 共11题; 题目175分, 表达5分, 合计180分.
\begin{exercise}原卷第1题 (30分). 定义 (i) 群; (ii) 群自同构; (iii) 二面体群; (iv) 理想; (v) PID; (vi) UFD.\end{exercise}
\begin{solution}
(i) 群的运算结合, 有双侧单位元及每个元素的双侧逆元. (ii) 自同构是群到自身的双射同态. (iii) $D_n$ 为正 $n$ 边形的全对称群, 本书阶约定为 $2n$; 可用 $r^n=s^2=e,srs=r^{-1}$ 描述. (iv) 理想是环的加法子群, 对任意环元素左右乘法吸收; 交换环只需一侧. (v) PID 为每个理想主的整环. (vi) UFD 为每个非零非单位都可分解成不可约元且仅差次序与相伴的整环.
\end{solution}
\begin{exercise}原卷第2题 (15分). (i) 证明共轭关系是等价关系. (ii) 列出 $S_5$ 的等价类.\end{exercise}
\begin{solution}
用单位元给自反, 共轭元的逆给对称, 两次共轭用乘积给传递. $S_5$ 的共轭类按5的分拆列出, 类型为 $1^5,2\,1^3,2^2\,1,3\,1^2,3\,2,4\,1,5$. 其大小分别为 $1,10,15,20,20,30,24$, 和为120. 大小公式 $5!/\prod_r(r^{m_r}m_r!)$ 来自循环内起点 $r$ 种重复与相同长度循环的 $m_r!$ 种重排.
\end{solution}
\begin{exercise}\label{ex:smallgroups}原卷第3题 (15分). 分类所有阶至多10的群.\end{exercise}
\begin{solution}
各阶的全部同构类型如下: 1阶平凡群; 2,3,5,7阶分别唯一 $C_2,C_3,C_5,C_7$; 4阶 $C_4,C_2^2$; 6阶 $C_6,S_3$; 8阶 $C_8,C_4\times C_2,C_2^3,D_4,Q_8$; 9阶 $C_9,C_3^2$; 10阶 $C_{10},D_5$.

素数阶与平方阶已由拉格朗日及第4章给出. 对 $2q$ ($q=3,5$), 唯一 Sylow $q$ 子群 $\langle a\rangle$ 正规, 柯西定理给阶2元素 $b$ 在其外. 群中每个元素为 $a^i$ 或 $a^ib$, 且 $bab^{-1}=a^u$, $u^2=1$ 模 $q$, 故 $u=1$ 或 $-1$. 前者循环直积, 后者二面体; $q=3$ 的二面体即 $S_3$.

8阶交换群由有限生成交换群结构定理给三类. 若非交换, 不能有阶8元素. 若所有非单位元阶2, 由 $(ab)^2=e$ 推得 $ab=ba$, 矛盾. 故取阶4元素 $a$, 子群指数2而正规, 选 $b$ 在其外. 共轭作用非平凡, 必 $bab^{-1}=a^{-1}$; 否则 $a,b$ 交换使群交换. $b^2\in\langle a\rangle$, 又与 $b$ 交换, 因此 $b^2=a^k$ 的 $k$ 必为0或2. $b^2=e$ 得 $D_4$, $b^2=a^2$ 得四元数群 $Q_8=\{\pm1,\pm i,\pm j,\pm k\}$, 其中 $i^2=j^2=k^2=ijk=-1$. 两个表示各最多8个规范词, 具体群已有8个元素, 因而无遗漏. $D_4$ 有5个阶2元素, $Q_8$ 只有1个, 故不同构.
\end{solution}
\begin{exercise}原卷第4题 (15分). (i) 陈述第二同构定理. (ii) 证明它.\end{exercise}
\begin{solution}
若 $H\le G,N\trianglelefteq G$, 则 $H/(H\cap N)\cong HN/N$. $HN$ 由正规性封闭; $H\to HN/N$, $h\mapsto hN$, 满射且核为 $H\cap N$, 第一同构定理给结论. 原课程对“第二”编号采用此版本, 不与第三同构定理混淆.
\end{solution}
\begin{exercise}原卷第5题 (15分). (i) 陈述 Sylow 定理. (ii) 若 $|G|=pqr$, 三素数互异, 证明 $G$ 不单.\end{exercise}
\begin{solution}
(i) 见习题\ref{ex:sylowstate}. (ii) 设 $p<q<r$ 且群单. $n_r\mid pq$, 模 $r$ 余1且大于1, 不可能为 $p$ 或 $q$, 只能 $pq$. 各阶 $r$ 子群除单位外不交, 共 $pq(r-1)$ 个元素. 又 $n_q$ 是 $pr$ 的因子, 大于1且模 $q$ 余1, 不可能为 $p$, 所以 $n_q\ge r$. 它们至少给 $r(q-1)$ 个不同阶 $q$ 元素, 与前者不交. 但
\[
pq(r-1)+r(q-1)=pqr+\bigl(r(q-1)-pq\bigr)>pqr,
\]
因 $r>q$ 且 $p\le q-1$ 得 $r(q-1)>q(q-1)\ge pq$. 矛盾.
\end{solution}
\begin{exercise}原卷第6题 (15分). 设 $R$ 交换, $I,J$ 为理想. (i) 素理想 $P\supseteq IJ$ 时证明 $P\supseteq I$ 或 $P\supseteq J$. (ii) 建立 $R/IJ$ 与 $R/(I\cap J)$ 的素理想自然双射. (iii) 给 $IJ\ne I\cap J$ 的例子.\end{exercise}
\begin{solution}
(i) 若 $I\not\subseteq P$, 取 $i\in I\setminus P$, 则对所有 $j\in J$, $ij\in P$, 素性迫使 $j\in P$. (ii) $IJ\subseteq I\cap J$. 包含 $IJ$ 的素理想由(i)包含一方, 因而包含交; 包含交当然包含积. 两商环的素理想因此对应同一批 $R$ 的素理想, 用投影原像与像得到双射. (iii) $R=\Z,I=J=2\Z$, 则积 $4\Z$, 交 $2\Z$.
\end{solution}
\begin{exercise}原卷第7题 (10分). 非域 UFD 是否必有无限多个两两不相伴的不可约元?\end{exercise}
\begin{solution}
否. 取 $R=\{a/b\in\Q:b\text{ 为奇数}\}=\Z_{(2)}$. 非零元素唯一写为 $2^ku$, $k\ge0$, $u$ 为单位 (分子分母皆奇). 非单位正是 $k>0$, 不可约元正是 $k=1$, 全与2相伴. 分解存在且唯一, 所以为非域 UFD, 却只有一个不可约元相伴类.
\end{solution}
\begin{exercise}原卷第8题 (10分). 举一个每个非零非单位可分解为不可约元却非 UFD 的整环.\end{exercise}
\begin{solution}
取 $\Z[\sqrt{-5}]$. 范数为正整数且每个非单位范数大于1, 严格下降归纳保证有限分解. $6=2\cdot3=(1+\sqrt{-5})(1-\sqrt{-5})$ 给不相伴的两种不可约分解; 具体不可约性见第6章的完整范数检验.
\end{solution}
\begin{exercise}原卷第9题 (15分). (i) 证明 $\Z[i]$ 欧几里得. (ii) $6-i$ 是否为素元?\end{exercise}
\begin{solution}
(i) 用范数 $a^2+b^2$; 对 $z/w$ 实虚部各选最近整数 $q$, 余项 $r=z-qw$ 满足 $N(r)\le N(w)/2<N(w)$. (ii) 范数 $37$ 为有理素数. 若 $6-i=ab$, 范数乘法性迫使一个因子范数1, 即单位, 故不可约. 欧几里得整环是 PID, 不可约即素元, 所以是.
\end{solution}
\begin{exercise}原卷第10题 (10分). 写出 $\F_5[x]$ 的全部不可约二次多项式.\end{exercise}
\begin{solution}
先列首一式. 对 $x^2+bx+c$, 完全平方表明无根当且仅当判别式 $b^2-4c$ 为非平方, 即2或3. 全部十个首一式为
\begin{align*}
&x^2+2,\ x^2+3,\ x^2+x+1,\ x^2+x+2,\\
&x^2+2x+3,\ x^2+2x+4,\ x^2+3x+3,\ x^2+3x+4,\\
&x^2+4x+1,\ x^2+4x+2.
\end{align*}
原题未限定首一, 所以还须取上述每式的任一非零常数倍, 共40个二次不可约多项式.
\end{solution}
\begin{exercise}原卷第11题 (25分). (i) 陈述高斯引理与艾森斯坦判别. (ii) 证明 $1+x^3+x^6$ 在 $\Q[x]$ 不可约. (iii) 证明 $1-t^2+t^5$ 在 $\Q[t]$ 不可约.\end{exercise}
\begin{solution}
(i) 本原多项式之积本原, 本原整数多项式在 $\Z[x]$ 与 $\Q[x]$ 的不可约性等价. 艾森斯坦条件为: 正次数整数多项式存在素数 $p$ 整除全部非最高次系数, $p$ 不整除最高系数, 且 $p^2$ 不整除常数项; 则该多项式在 $\Q[x]$ 上不可约. 完整证明见引理\ref{lem:gauss}及定理\ref{thm:eisenstein}. (ii) 令 $x=y+1$, 得
\[
y^6+6y^5+15y^4+21y^3+18y^2+9y+3.
\]
素数3整除所有非最高系数, 不整除最高系数且9不整除常数3, 判别成立; 平移可逆, 故原式不可约.

(iii) 模2为 $h=t^5+t^2+1$, 在0,1无根. 若五次式可约且无一次因子, 必有不可约二次因子, 因可能次数分拆至少一项不大于2. $\F_2$ 唯一不可约二次为 $t^2+t+1$. 模此式 $t^2=t+1$, $t^3=1$, $t^5=t^2$, 故 $h\equiv1$, 不整除. 因而 $h$ 不可约, 约化判别给有理不可约性.
\end{solution}

\originalchapter{来源对应与范围边界}
23份教师 PDF 共149个实际页面; 第24讲为复习, 官网没有讲义. 各文件页数、原始直链、SHA256及许可检查见 source-manifest.json, 这些原件不在中文源码包中. 对应关系如下; “择要”明确表示未覆盖整份讲义, 不能据此宣称逐页全译.
\begin{longtable}{@{}p{.12\textwidth}p{.32\textwidth}p{.49\textwidth}@{}}
\toprule 原讲次&原主题&中文册与边界\\\midrule\endhead
1--4&群、子群、陪集、循环群&第1--2章, 完整核心论证, 合并重复例子\\
5--6&置换与共轭&第3章, 完整核心论证\\
7--10&同构、同态、商及同构定理&第2章, 完整核心论证\\
11&交错群&第3章符号, 作业6中补完整单性证明\\
12&群表示与小阶群&第3章二面体表示, 模拟期末第3题分类至10阶; 一般自由群构造未展开\\
13&Sylow&第4章完整证明; 作业和模拟卷完整公开题面题解\\
14--16&环、基本性质、同态与理想&第5章, 完整核心论证\\
17--20&分式域、素/极大理想、特殊及欧几里得整环&第5--6章, 完整核心论证\\
21&多项式&第7章及作业11, 完整核心论证\\
22&快速素性检验&第14章择要: Fermat、多项式判别、平方根见证; 不含 AKS 参数界与完整复杂度证明, 不声称完整覆盖\\
23&群作用与自同构&群作用在第4章, 域自同构第11--13章, 多项式自同构在作业11\\
\bottomrule
\end{longtable}
RES.18-011第1--6讲的基础群论并入第1--3章; 第17--19、21--23讲的作用、稳定子、共轭与 Sylow 并入第3--4章. 其第7--11讲线性代数仅作为模论应用的先修参考, 不逐讲重写. 第12--16、20、24--35讲的正交几何、离散群、二十面体、双线性/厄米空间、谱定理、线性Lie群及 Hilbert 第三问题未纳入.

RES.18-012第8--13讲的环、理想、分解并入第5--7章; 第19--22讲的模、关系矩阵、Smith 与分解构成第8--9章; 第24--27、29--31讲的域、扩张、有限域、Galois 构成第10--13章. 第23讲 Noether 环仅使用 PID 升链停止的自足特例, 未证明一般 Hilbert 基定理. 第1--7讲及附录的群表示与特征标、第14--18讲数域理想分解与类群、第28讲函数域几何、第32--35讲一般根式求解与判别式理论未完整纳入. 本原元、Artin 定理及各核心定理中的省略步骤由本册补写; 不是英文整页替换为中文标题.
原18.703第22讲关于模奇合数的平方根1数量的引理漏掉素数幂情形, 第14章以模9反例和正确条件修正. 原作业8把环直积当作一般环余积的附注不成立, 第5章给反例. 作业2第3题需排除平凡群并澄清“没有非平凡真子群”; 作业6第12题的交错群单性陈述需 $n\ge3$. 作业11在文件标题中误写作业9, 本册按官网作业顺序编号11并保留映射. 这些改动是数学勘正, 不宣称官方已勘误.

同源合订PDF与单讲PDF内容重复, 学生原 TeX 包与合订PDF是同一笔记的不同格式, 故不另算学科或重复抄入正文. 英文官方包只作来源归档, 原图不复用; 本册唯一域格图由 TikZ 独立绘制.

\begin{thebibliography}{9}
\bibitem{mck} James McKernan. \textit{18.703 Modern Algebra, Spring 2013}. MIT OpenCourseWare. \href{https://ocw.mit.edu/courses/18-703-modern-algebra-spring-2013/pages/lecture-notes/}{23份教师讲义目录}.
\bibitem{alg1} MIT OpenCourseWare. \textit{RES.18-011 Algebra I Student Notes, Fall 2021}. \href{https://ocw.mit.edu/courses/res-18-011-algebra-i-student-notes-fall-2021/pages/student-notes/}{学生笔记与官方源码}. 未经教师核准.
\bibitem{alg2} MIT OpenCourseWare. \textit{RES.18-012 Algebra II Student Notes, Spring 2022}. \href{https://ocw.mit.edu/courses/res-18-012-algebra-ii-student-notes-spring-2022/pages/student-notes/}{学生笔记与官方源码}. 未经教师核准.
\bibitem{terms} MIT OpenCourseWare. \href{https://ocw.mit.edu/pages/privacy-and-terms-of-use/}{使用条款}; Creative Commons, \href{https://creativecommons.org/licenses/by-nc-sa/4.0/}{CC BY-NC-SA 4.0}.
\end{thebibliography}
\end{document}
